% !TEX program = lualatex
% ---------------------------------------------------------------------------------------------
%  Readback note 10: Vanishing of theta descents: reflection sums at the level, a maximum principle for Hecke polynomials, and adding hyperbolic planes
%
%  Build:  python3 tools/rb.py note <slug>
%
%  One unit of a Lean artifact, rewritten as ordinary mathematics. The prose is hand-written. Lean names,
%  line ranges and counts come from data.tex; quoted Lean statements are copied from the artifact by the build.
%
%  Rules of the text: standard terminology; every symbol is explained where it first appears and listed again
%  in the notation at the end; no letter has two meanings; a proof is complete or says what it takes for granted;
%  Lean names appear only in the LEAN lines and in the concordance.
%
%  The signature rule (the build enforces it): a function, relation or predicate that the artifact defines is
%  written only through its notation macro, which the build makes from the spec (`notation`) with one slot for
%  each explicit argument of the Lean definition. No argument is fixed and dropped, also not "for brevity" and
%  not for one section. Every argument and every hypothesis of a stated declaration stands inside its
%  statement, or in a Setting that the statement cites by number; only implicit and instance arguments (the
%  ambient objects) may stand in the conventions.
% ---------------------------------------------------------------------------------------------
\documentclass[11pt,a4paper]{article}
\ifdefined\pdfvariable\pdfvariable suppressoptionalinfo \numexpr 1+2+4+8\relax\fi
\usepackage[a4paper,hmargin=35mm,top=22mm,bottom=27mm,footskip=12mm,nohead]{geometry}

\usepackage{fontspec}
\directlua{
  luaotfload.add_fallback("gfallback", {"STIXTwoMath-Regular.otf:mode=node;", "DejaVuSans.ttf:mode=node;"})
  luaotfload.add_fallback("gmonofallback", {"DejaVuSansMono.ttf:mode=node;", "STIXTwoMath-Regular.otf:mode=node;"})
}
\setmainfont{NewCM10}[Extension=.otf,UprightFont=*-Book,ItalicFont=*-BookItalic,
  BoldFont=*-Bold,BoldItalicFont=*-BoldItalic,SmallCapsFont=*-Book,SmallCapsFeatures={Letters=SmallCaps},
  RawFeature={fallback=gfallback}]
\setsansfont{NewCMSans10}[Extension=.otf,UprightFont=*-Book,ItalicFont=*-BookOblique,
  BoldFont=*-Bold,BoldItalicFont=*-BoldOblique,RawFeature={fallback=gfallback}]
\setmonofont{JetBrainsMono}[Extension=.otf,UprightFont=*-Regular,ItalicFont=*-Italic,
  BoldFont=*-SemiBold,BoldItalicFont=*-SemiBoldItalic,Scale=MatchLowercase,
  RawFeature={-calt;-liga;fallback=gmonofallback}]
\usepackage{amsmath}
\usepackage{amsthm}
\usepackage[math-style=TeX]{unicode-math}
\setmathfont{NewCMMath-Book.otf}
\emergencystretch=2.5em\relax

\usepackage{xcolor}
\usepackage{tikz}
\usetikzlibrary{calc,arrows.meta,positioning,backgrounds}
\usepackage[protrusion=true,expansion=true]{microtype}
\usepackage{enumitem}
\usepackage{needspace}
\usepackage{booktabs}
\usepackage{array}
\usepackage{longtable}
\usepackage{fancyhdr}
\usepackage{lastpage}
\usepackage{multicol}
\usepackage[hidelinks,pdfusetitle]{hyperref}

\definecolor{ink}{HTML}{0F172A}
\definecolor{inksoft}{HTML}{475569}
\definecolor{inkfaint}{HTML}{94A3B8}
\definecolor{rulegrey}{HTML}{CBD5E1}
\definecolor{arith}{HTML}{C2410C}
\definecolor{est}{HTML}{7C3AED}
\definecolor{ana}{HTML}{0891B2}
\definecolor{ok}{HTML}{047857}
\definecolor{leancol}{HTML}{3730A3}

\title{Vanishing of theta descents: reflection sums at the level, a maximum principle for Hecke polynomials, and adding hyperbolic planes}
\author{grh-compact}
\date{}

% ------------------------------------------------------------------------------------- macros
\newcommand{\leanfont}{\ttfamily\fontsize{8.6}{10.5}\selectfont\color{leancol}}
\newcommand{\lean}[1]{\csname LN@#1\endcsname}
\newcommand{\fillin}[1]{{\color{arith}\itshape[#1]}}  % a place of the template that is not written yet; the build reports every one that is left
\newcommand{\leanquote}[1]{\csname LQ@#1\endcsname}
% generated by rb.py note; do not edit
\newcommand{\ConcDecls}{8}
\newcommand{\Pin}{72fead10}
\newcommand{\PinText}{that of commit {\leanfont 72fead10}}
\newcommand{\ArtifactName}{grh-compact}
\expandafter\newcommand\csname LN@existsatPlaceeq\endcsname{\href{https://github.com/lsernau/grh-lean/blob/72fead102f663d5c7a407964a13dc0cf9d977207/grh-compact/src/New/SiegelWeil/AuxPairVanishing.lean\#L12-L15}{{\leanfont AuxPair.\allowbreak{}exists\_atPlace\_eq}}}
\expandafter\newcommand\csname LN@existsannihilatororcongrinclMatrix\endcsname{\href{https://github.com/lsernau/grh-lean/blob/72fead102f663d5c7a407964a13dc0cf9d977207/grh-compact/src/New/SiegelWeil/AuxPairVanishing.lean\#L20-L33}{{\leanfont AuxPair.\allowbreak{}exists\_annihilator\_or\_congr\_inclMatrix}}}
\expandafter\newcommand\csname LN@smulchosenGLRepeq\endcsname{\href{https://github.com/lsernau/grh-lean/blob/72fead102f663d5c7a407964a13dc0cf9d977207/grh-compact/src/New/SiegelWeil/AuxPairVanishing.lean\#L37-L49}{{\leanfont AuxPair.\allowbreak{}smul\_chosenGLRep\_eq}}}
\expandafter\newcommand\csname LN@isAutomorphicstaralgFormCoeff\endcsname{\href{https://github.com/lsernau/grh-lean/blob/72fead102f663d5c7a407964a13dc0cf9d977207/grh-compact/src/New/SiegelWeil/AuxPairVanishing.lean\#L51-L52}{{\leanfont AuxPair.\allowbreak{}isAutomorphic\_star\_algFormCoeff}}}
\expandafter\newcommand\csname LN@descentsVanishofreflectionSum\endcsname{\href{https://github.com/lsernau/grh-lean/blob/72fead102f663d5c7a407964a13dc0cf9d977207/grh-compact/src/New/SiegelWeil/AuxPairVanishing.lean\#L59-L158}{{\leanfont AuxPair.\allowbreak{}descentsVanish\_of\_reflectionSum}}}
\expandafter\newcommand\csname LQ@descentsVanishofreflectionSum\endcsname{\leanquoteline{0}{lemma AuxPair.descentsVanish\_of\_reflectionSum [WeilData E]}\leanquoteline{4}{\{y\} u (hQ : ∀ i, O.ψ (O.twistPlace i) ≠ 1)}\leanquoteline{4}{(hM : ∀ i (v : ∀ \_, \_), v i ≠ 0 → O.formAction.asAlgebraHom (piTensorElem fun i ↦}\leanquoteline{2}{reflectionSum (O.ψ \_) (v i)) y = 0) (h0 : ∀ P, lift (pairThetaKer \_ \_ O.a 1).swap (integrated \_}\leanquoteline{2}{(polyVector \_ P) O.auxDatum.1) (star (algFormCoeff 1 u y.1)) = 0) : DescentsVanish O.W (stdSpace E}\leanquoteline{2}{n) O.a O.places.conjClosureᶜ (star (algFormCoeff 1 u y.1))}}
\expandafter\newcommand\csname LN@eqzeroofheckePolyOpeqzero\endcsname{\href{https://github.com/lsernau/grh-lean/blob/72fead102f663d5c7a407964a13dc0cf9d977207/grh-compact/src/New/SiegelWeil/DescentsVanishPlanes.lean\#L17-L28}{{\leanfont eq\_zero\_of\_heckePolyOp\_eq\_zero}}}
\expandafter\newcommand\csname LN@integralsourceTermeqzero\endcsname{\href{https://github.com/lsernau/grh-lean/blob/72fead102f663d5c7a407964a13dc0cf9d977207/grh-compact/src/New/SiegelWeil/DescentsVanishPlanes.lean\#L35-L44}{{\leanfont DescentsVanish.\allowbreak{}integral\_sourceTerm\_eq\_zero}}}
\expandafter\newcommand\csname LN@descentsVanishaddplanes\endcsname{\href{https://github.com/lsernau/grh-lean/blob/72fead102f663d5c7a407964a13dc0cf9d977207/grh-compact/src/New/SiegelWeil/DescentsVanishPlanes.lean\#L49-L109}{{\leanfont descentsVanish\_add\_planes}}}
\expandafter\newcommand\csname LQ@descentsVanishaddplanes\endcsname{\leanquoteline{0}{lemma descentsVanish\_add\_planes [IsBaseAdapted D] (hY : V.IsAutomorphic Y) \{w\} (hw :}\leanquoteline{2}{w.IsSplit) (ha : CharUnramifiedAt a.baseChange.finPart (star w))}\leanquoteline{2}{(hw' : star w ∈ S) (hK : Y ∈ sphericalAt (finRightRep ..) w) \{X\} (hX : ‖X‖ = (√w.q \textasciicircum{} card r *}\leanquoteline{2}{w.q \textasciicircum{} (card n - card r - 1))⁻¹) (hP : heckePolyOp (finRightRep ..) hw X Y = 0) (h : 2 * card}\leanquoteline{2}{r + 2 * card p + 3 < card n) : DescentsVanish W' V (a * upsilonU1 \textasciicircum{} card p) S Y}}
\newcommand{\ConcordanceRows}{\crow{existsatPlaceeq}{\href{https://github.com/lsernau/grh-lean/blob/72fead102f663d5c7a407964a13dc0cf9d977207/grh-compact/src/New/SiegelWeil/AuxPairVanishing.lean\#L12-L15}{{\leanfont AuxPair.\allowbreak{}exists\_atPlace\_eq}}}{AuxPairVanishing.lean}{12--15}
\crow{existsannihilatororcongrinclMatrix}{\href{https://github.com/lsernau/grh-lean/blob/72fead102f663d5c7a407964a13dc0cf9d977207/grh-compact/src/New/SiegelWeil/AuxPairVanishing.lean\#L20-L33}{{\leanfont AuxPair.\allowbreak{}exists\_annihilator\_or\_congr\_inclMatrix}}}{AuxPairVanishing.lean}{20--33}
\crow{smulchosenGLRepeq}{\href{https://github.com/lsernau/grh-lean/blob/72fead102f663d5c7a407964a13dc0cf9d977207/grh-compact/src/New/SiegelWeil/AuxPairVanishing.lean\#L37-L49}{{\leanfont AuxPair.\allowbreak{}smul\_chosenGLRep\_eq}}}{AuxPairVanishing.lean}{37--49}
\crow{isAutomorphicstaralgFormCoeff}{\href{https://github.com/lsernau/grh-lean/blob/72fead102f663d5c7a407964a13dc0cf9d977207/grh-compact/src/New/SiegelWeil/AuxPairVanishing.lean\#L51-L52}{{\leanfont AuxPair.\allowbreak{}isAutomorphic\_star\_algFormCoeff}}}{AuxPairVanishing.lean}{51--52}
\crow{descentsVanishofreflectionSum}{\href{https://github.com/lsernau/grh-lean/blob/72fead102f663d5c7a407964a13dc0cf9d977207/grh-compact/src/New/SiegelWeil/AuxPairVanishing.lean\#L59-L158}{{\leanfont AuxPair.\allowbreak{}descentsVanish\_of\_reflectionSum}}}{AuxPairVanishing.lean}{59--158}
\crow{eqzeroofheckePolyOpeqzero}{\href{https://github.com/lsernau/grh-lean/blob/72fead102f663d5c7a407964a13dc0cf9d977207/grh-compact/src/New/SiegelWeil/DescentsVanishPlanes.lean\#L17-L28}{{\leanfont eq\_zero\_of\_heckePolyOp\_eq\_zero}}}{DescentsVanishPlanes.lean}{17--28}
\crow{integralsourceTermeqzero}{\href{https://github.com/lsernau/grh-lean/blob/72fead102f663d5c7a407964a13dc0cf9d977207/grh-compact/src/New/SiegelWeil/DescentsVanishPlanes.lean\#L35-L44}{{\leanfont DescentsVanish.\allowbreak{}integral\_sourceTerm\_eq\_zero}}}{DescentsVanishPlanes.lean}{35--44}
\crow{descentsVanishaddplanes}{\href{https://github.com/lsernau/grh-lean/blob/72fead102f663d5c7a407964a13dc0cf9d977207/grh-compact/src/New/SiegelWeil/DescentsVanishPlanes.lean\#L49-L109}{{\leanfont descentsVanish\_add\_planes}}}{DescentsVanishPlanes.lean}{49--109}}
% the notations of the defined functions: one slot for each explicit argument of the Lean declaration
\newcommand{\UnitMembers}{8}
\newcommand{\UnitInNote}{8}
\newcommand{\UnitProofs}{7}
\newcommand{\UnitSteps}{89}
\newcommand{\UnitAllSteps}{104}


\linespread{1.08}                       % 11 pt on 14.7 pt
\setlength{\parindent}{0pt}
\setlength{\parskip}{5.5pt plus 1.5pt minus 0.5pt}
\setlength{\abovedisplayskip}{7.5pt plus 2pt minus 3pt}
\setlength{\belowdisplayskip}{7.5pt plus 2pt minus 3pt}
\setlength{\abovedisplayshortskip}{3pt plus 2pt}
\setlength{\belowdisplayshortskip}{6pt plus 2pt minus 2pt}
\emergencystretch=1.5em
\clubpenalty=10000 \widowpenalty=10000 \displaywidowpenalty=10000
\relpenalty=7000 \binoppenalty=9000
\raggedbottom
\pagestyle{fancy}\fancyhf{}\renewcommand{\headrulewidth}{0pt}
\fancyfoot[C]{\sffamily\fontsize{7.6}{9}\selectfont\color{inkfaint}\thepage\,/\,\pageref*{LastPage}}
\setlist{topsep=2pt,partopsep=0pt,parsep=0pt,itemsep=4pt,leftmargin=2.1em,labelsep=0.55em}

\newcommand{\Q}{\mathbb{Q}}
\newcommand{\Z}{\mathbb{Z}}
\newcommand{\F}{\mathbb{F}}
\newcommand{\len}{\operatorname{len}}
\newcommand{\cok}{\operatorname{cok}}
\newcommand{\End}{\operatorname{End}}
\newcommand{\Hom}{\operatorname{Hom}}
\newcommand{\Gal}{\operatorname{Gal}}
\newcommand{\diag}{\operatorname{diag}}
\newcommand{\kw}[1]{\emph{#1}}
\newcommand{\m}[1]{\mbox{$#1$}}        % a short formula that must not be broken across lines
\newcommand{\head}[1]{\par\addvspace{17pt plus 4pt minus 2pt}\Needspace{6\baselineskip}%
  \noindent{\fontsize{13.2}{16}\selectfont\bfseries\color{ink}#1}\par\nopagebreak\vspace{-1pt}}
\newcommand{\sect}[2]{\refstepcounter{section}\label{#2}\par\addvspace{19pt plus 4pt minus 2pt}\Needspace{7\baselineskip}%
  \noindent\llap{{\fontsize{13.2}{16}\selectfont\bfseries\color{inkfaint}\thesection}\hspace{9pt}}%
  {\fontsize{13.2}{16}\selectfont\bfseries\color{ink}#1}\par\nopagebreak\vspace{-1pt}}
\newcommand{\subhead}[1]{\par\addvspace{9pt plus 2pt}\Needspace{4\baselineskip}\noindent\textbf{#1}\enspace\ignorespaces}
% where a statement lives in the Lean archive
\newcommand{\leantag}{{\sffamily\bfseries\fontsize{6.3}{8}\selectfont\color{leancol}\addfontfeatures{LetterSpace=7}LEAN}}
\newcommand{\sep}{\unskip\nobreak\hspace{5pt}{\color{inkfaint}\textperiodcentered}\hspace{5pt plus 2pt}\ignorespaces}
\newcommand{\leanline}[1]{\par\nopagebreak\vspace{-2.5pt}{\fontsize{8.9}{12.4}\selectfont\color{inksoft}\raggedright
  \noindent\llap{\leantag\hspace{9pt}}#1\par}\vspace{1pt}}
% a displayed statement in a tinted box: #1 colour, #2 text
\newcommand{\statement}[2]{\par\addvspace{7pt}\begingroup\setlength{\fboxsep}{8.5pt}\setlength{\fboxrule}{0.6pt}%
  \noindent\fcolorbox{#1!30!white}{#1!3!white}{\parbox{\dimexpr\linewidth-2\fboxsep-2\fboxrule\relax}{#2}}%
  \endgroup\par\addvspace{7pt}}
\newcommand{\nt}[2]{\par\noindent\hangindent=1.1em\hangafter=1 #1\enspace{\color{inksoft}#2}\par\vspace{1pt}}
\newcommand{\legdot}[1]{\tikz[baseline=-0.62ex]\fill[#1] (0,0) circle (2.5pt);\hspace{3pt}}
\newcommand{\crow}[4]{\csname CS@#1\endcsname & #2 & \textcolor{inksoft}{#3} & \textcolor{inksoft}{#4}\\}
\newcommand{\leanquoteline}[2]{\par\noindent\hangindent=\dimexpr#1\fontdimen2\font+2.2em\relax\hangafter=1
  \hspace*{#1\fontdimen2\font}#2}
\newcommand{\quoted}[1]{\par\vspace{3pt}{\ttfamily\fontsize{7.5}{10.4}\selectfont\color{leancol}\raggedright
  \setlength{\parskip}{0pt}\leanquote{#1}\par}\vspace{4pt}}
\newcommand{\bx}[3]{{\sffamily\bfseries\fontsize{7.3}{9.2}\selectfont\color{#1!82!black}#2}\\[1.5pt]#3}
\newcommand{\cs}[2]{\expandafter\newcommand\csname CS@#1\endcsname{#2}}

% theorems
\newtheoremstyle{plainx}{7pt plus 2pt}{7pt plus 2pt}{\itshape}{}{\bfseries}{.}{0.5em}{}
\newtheoremstyle{uprightx}{7pt plus 2pt}{7pt plus 2pt}{}{}{\bfseries}{.}{0.5em}{}
\theoremstyle{plainx}
\newtheorem{theorem}{Theorem}[section]
\newtheorem{lemma}[theorem]{Lemma}
\newtheorem{proposition}[theorem]{Proposition}
\newtheorem{corollary}[theorem]{Corollary}
\theoremstyle{uprightx}
\newtheorem{definition}[theorem]{Definition}
\newtheorem{setting}[theorem]{Setting}          % hypotheses shared by several statements; each of them cites it by number
\newtheorem{remark}[theorem]{Remark}
\newtheorem{example}[theorem]{Example}
\renewcommand{\qedsymbol}{{\setlength{\fboxsep}{0pt}\setlength{\fboxrule}{0.5pt}\fbox{\rule{0pt}{0.46em}\hspace{0.46em}}}}
% what holds throughout the note: the ambient objects only (implicit and instance arguments), never a hypothesis
\newenvironment{conventions}{\par\addvspace{5pt}\noindent\textbf{Conventions.}\enspace\ignorespaces}{\par\addvspace{3pt}}
\makeatletter
\renewenvironment{proof}[1][Proof]{\par\pushQED{\qed}\normalfont\topsep0pt\partopsep0pt
  \trivlist\item[\hskip\labelsep\itshape #1.]\ignorespaces}{\popQED\endtrivlist\@endpefalse}
\makeatother

% where each declaration of the concordance is treated: \cs{key}{place}, the place naming a \label or a lettered item
\cs{existsatPlaceeq}{Lemma~\ref{l:atplace}}
\cs{existsannihilatororcongrinclMatrix}{Lemma~\ref{l:annihilator}}
\cs{smulchosenGLRepeq}{Lemma~\ref{l:level}}
\cs{isAutomorphicstaralgFormCoeff}{Lemma~\ref{l:automorphic}}
\cs{descentsVanishofreflectionSum}{Theorem 1}
\cs{eqzeroofheckePolyOpeqzero}{Proposition~\ref{p:max}}
\cs{integralsourceTermeqzero}{Lemma~\ref{l:zero}}
\cs{descentsVanishaddplanes}{Theorem 2}
% notation of this note; objects of other units and of the background library are written with all their arguments
\newcommand{\CC}{\mathbb{C}}
\newcommand{\AAf}{\mathbb{A}_{\mathrm{f}}}
\newcommand{\FF}{\mathbb{F}}
\newcommand{\oo}{\mathfrak{o}}
\newcommand{\ohat}{\hat{\mathfrak{o}}}
\newcommand{\Ginf}[1]{G_{\infty}(#1)}
\newcommand{\Gfin}[1]{G_{\mathrm{f}}(#1)}
\newcommand{\Gad}[1]{G(#1)}
\newcommand{\Kz}[1]{\mathbf{K}_{0}(#1)}
\newcommand{\std}[2]{\mathrm{I}_{#2}(#1)}
\newcommand{\one}[1]{1_{#1}}
\newcommand{\WO}[1]{W(#1)}
\newcommand{\sgO}[1]{\varepsilon_{#1}}
\newcommand{\incO}[1]{\operatorname{inc}_{#1}}
\newcommand{\IncM}[1]{\operatorname{Inc}_{#1}}
\newcommand{\kap}[1]{\kappa_{#1}}
\newcommand{\eO}[1]{e_{#1}}
\newcommand{\psiO}[2]{\psi_{#1,#2}}
\newcommand{\MO}[1]{M_{#1}}
\newcommand{\QO}[1]{Q_{#1}}
\newcommand{\plO}[1]{\operatorname{pl}(#1)}
\newcommand{\SigO}[1]{\Sigma_{#1}}
\newcommand{\ChGL}[2]{\mathrm{GL}^{\mathrm{ch}}_{#2}(#1)}
\newcommand{\chO}[1]{\chi_{#1}}
\newcommand{\LO}[1]{\mathbf{L}(#1)}
\newcommand{\KO}[1]{\mathbf{K}(#1)}
\newcommand{\etaO}[1]{\eta_{#1}}
\newcommand{\tetaO}[1]{\tilde{\eta}_{#1}}
\newcommand{\Ups}[1]{\Upsilon(#1)}
\newcommand{\PsiO}[2]{\Psi_{#1}(#2)}
\newcommand{\FormO}[1]{\operatorname{Form}(#1)}
\newcommand{\Vk}[2]{\mathbb{V}_{#1,#2}}
\newcommand{\rhoW}[2]{\rho_{#1,#2}}
\newcommand{\Desc}[4]{\operatorname{Desc}^{#1,#2}_{#3,#4}}
\newcommand{\pvO}[2]{v_{#1,#2}}
\newcommand{\red}[1]{\operatorname{red}_{#1}}
\newcommand{\coef}[2]{#2_{#1}}
\newcommand{\chact}[3]{[#3]_{#1,#2}}
\newcommand{\Refl}[2]{\operatorname{Refl}(#1,#2)}
\newcommand{\PS}[1]{\mathrm{P}^{#1}}
\newcommand{\DV}[5]{\mathrm{DV}(#1,#2;#3;#4;#5)}
\newcommand{\Hk}[3]{H^{#1}_{#2}(#3)}
\newcommand{\Rt}[1]{\mathrm{R}_{#1}}
\newcommand{\ZT}[3]{\operatorname{Z}_{#1,#2}[#3]}
\newcommand{\hyp}[2]{#1\oplus\mathbb{H}^{#2}}
\newcommand{\cosI}[2]{\mathbf{1}_{#2+#1\operatorname{Box}}}
\newcommand{\vacf}{\varphi^{\circ}}
\newcommand{\thd}{\theta^{\Pi\prime}}
\newcommand{\lnm}[1]{{\leanfont #1}}
\newcommand{\conj}[1]{\overline{#1}}
\newcommand{\dd}{\mathrm{d}}
\newcommand{\ee}{\mathrm{e}}

\begin{document}
\color{ink}

{\raggedright
{\sffamily\bfseries\fontsize{7.2}{9}\selectfont\color{leancol}\addfontfeatures{LetterSpace=8}READBACK NOTE 10\par}
\vspace{5pt}
{\fontsize{21.5}{25.5}\selectfont\bfseries Vanishing of theta descents: reflection sums at the level, a maximum principle for Hecke polynomials, and adding hyperbolic planes\par}
\vspace{5pt}
{\fontsize{12.4}{16}\selectfont\itshape\color{inksoft} One unit of the Lean artifact \ArtifactName, read back into
mathematics\par}}
\vspace{8pt}
{\color{rulegrey}\hrule height 0.6pt}
\vspace{7pt}
\thispagestyle{empty}

\textbf{Summary.} A function~$Y$ on a definite unitary group has \emph{vanishing descents} to a space~$W$ if
its theta lifts to the unitary group of~$W$ are zero for all data that are unramified at a given set of places.
The unit proves two transfers of this property. First, for the auxiliary pair of the second note: if the
descents vanish for the single test datum at the level, with all polynomial vectors at the infinite places,
then they vanish for all data unramified outside the level, provided the form is killed by the reflection
sums at the chosen places and the character of the level is non-trivial at the twist places. Second, vanishing
climbs the theta tower: from~$W$ to~$W$ plus~$\nu$ hyperbolic planes, as long as \m{2r+2\nu+3<n} and $Y$ is
killed by a Hecke polynomial at one split place at a point that is too small to be a Hecke root of a lift. The
second transfer rests on a maximum principle: a bounded function killed by a Hecke polynomial with a small
argument is zero.

\statement{leancol}{\fontsize{9.6}{13.2}\selectfont\textbf{What this note is.} The artifact \ArtifactName{} is a
body of Lean code: the compact edition of a formal proof, due to L.~Alpöge with the AI system Claude, whose final
theorems are the Riemann hypothesis and its generalisation to Dirichlet \emph{L}-functions. That headline is a
claim. The authors report that every file compiles and that Lean's kernel accepted the proofs; an independent
build is recorded in the handout's record of checks when it reports. No mathematician has read the argument
through, and nothing here says that the Riemann hypothesis is proved. This note treats one unit of the artifact:
all \UnitInNote{} of its \UnitMembers{} declarations. The proofs are given in full relative to inputs that are
stated and not proved: the auxiliary pair and its objects (Setting~\ref{s:aux}, from the unit on the rows), and
results of the background library on the Weil representation, on Hecke operators at split places and on the
decay of theta kernels (the facts (Q1) to (Q11)). One step, the bound for the theta kernel near the cusps in the
proof of Theorem~2, is described through the library's estimate and was not traced in its constants
(§\ref{s:notes}, note~3). Every statement was compared with its Lean text, argument by argument. The Lean text
read is \PinText; Appendix~A is the concordance of \ConcDecls{} declarations. This is the third of three notes on
the theta correspondence behind Theorem~C; it uses the notation of the note \emph{The test datum of an auxiliary
pair, joint harmonics with staircase weights, and the theta lift as a conjugate descent}.}

\sect{The statements}{s:intro}
The objects are defined in §\ref{s:setting}. In brief: \m{\Desc{W}{V}{\chi}{1}(f,Y)} is the theta descent of a
function~$Y$ on the adelic group of a definite space~$V$ to the adelic group of~$W$, with datum~$f$, for the Weil
representation with characters $\chi$ and~$1$; \m{\DV{W}{V}{\chi}{S}{Y}} says that it vanishes for all data
unramified at the set of places~$S$; for an auxiliary pair~$O$, \m{\plO{O}} is the set of its chosen and twist
places, $\Ups{O}$ its finite datum, and \m{\coef{L}{\Xi}} the scalar function of a form \m{\Xi\in\FormO{O}} and a
linear form~$L$; \m{\Refl{\psi}{y}} is a reflection sum in the group algebra of a general linear group over a
finite field; and \m{\Hk{\pi}{w}{T}} is the Hecke polynomial at a split place~$w$, evaluated at a complex
number~$T$.

\statement{arith}{\phantomsection\label{t:1}\textbf{Theorem 1} (from the test datum to all data unramified
outside the level).\enspace \emph{In Setting~\ref{s:weil}, let $O$ be an auxiliary pair as in
Setting~\ref{s:aux}, \m{\Xi\in\FormO{O}} and $L$ a linear form on \m{\Vk{n}{\kap{O}}}. Suppose:}
\begin{enumerate}[label=\textup{(\alph*)},leftmargin=2.3em,itemsep=1pt,topsep=2pt]
\item \emph{\m{\psiO{O}{w}\ne1} for every twist place \m{w\in\QO{O}};}
\item \emph{for every chosen place \m{w_{0}\in\MO{O}} and every family \m{y=(y_{w})_{w\in\MO{O}}} of vectors
\m{y_{w}\in\FF_{w}^{r}} with \m{y_{w_{0}}\ne0}, the element \m{\bigotimes_{w\in\MO{O}}\Refl{\psiO{O}{w}}{y_{w}}} of the group
algebra of \m{\ChGL{\SigO{O}}{r}} annihilates~$\Xi$ in the action of Setting~\ref{s:aux}\textup{(c)};}
\item \emph{for every polynomial~$p$ in the variables \m{X_{\tau,i,c}},
\m{\Desc{\WO{O}}{\std{E}{n}}{\chO{O}}{1}\bigl(\pvO{O}{p}\otimes\Ups{O},\,\conj{\coef{L}{\Xi}}\bigr)=0}.}
\end{enumerate}
\emph{Then \m{\DV{\WO{O}}{\std{E}{n}}{\chO{O}}{S}{\conj{\coef{L}{\Xi}}}} holds for the set~$S$ of all finite places
of~$E$ that are neither in \m{\plO{O}} nor conjugate to a place of \m{\plO{O}}.}}
\leanline{\lean{descentsVanishofreflectionSum}}

\statement{arith}{\phantomsection\label{t:2}\textbf{Theorem 2} (climbing the theta tower).\enspace \emph{In
Setting~\ref{s:weil}, let $W$, $W'$ and~$V$ be hermitian spaces with points, of signs $\varepsilon$ on~$[r]$,
$\varepsilon'$ on~$[r']$ and \m{\one{n}} on~$[n]$, let $\chi$ be an automorphic character, $S$ a set of finite
places of~$E$ and $Y$ a complex function on \m{\Gad{\one{n}}}, and suppose \m{\DV{W}{V}{\chi}{S}{Y}}. Let $\Pi$ be an
adapted splitting of~$W'$ over~$W$ with planes indexed by a finite set~$[\nu]$ of $\nu$ elements
\textup{(Definition~\ref{d:split})}. Let $w$ be a finite place of~$E$ and $T$ a complex number. Suppose that $Y$
is automorphic on~$V$; that $w$ is split, \m{w\ne\conj{w}}; that the base change of~$\chi$ is unramified
at~\m{\conj{w}}; that \m{\conj{w}\in S}; that $Y$ is spherical at~$w$; that
\m{|T|=\bigl(\sqrt{q_{w}}^{\,r}\,q_{w}^{\,n-r-1}\bigr)^{-1}}; that \m{\Hk{\Rt{\one{n}}}{w}{T}\,Y=0}; and that
\m{2r+2\nu+3<n}. Then \m{\DV{W'}{V}{\chi\upsilon_{1}^{\nu}}{S}{Y}}.}}
\leanline{\lean{descentsVanishaddplanes}}

The space $W'$ is $W$ plus $\nu$ hyperbolic planes, and the spaces \m{\hyp{W}{\nu}}, \m{\nu=0,1,2,\dots}, form
the Witt tower of~$W$; the character changes by~$\upsilon_{1}$ with every plane. The number
\m{q_{w}^{-(n-1-r/2)}} is the absolute value of the smallest of the points at which the Hecke transfer (Q9) makes
the Hecke polynomial vanish on theta lifts from a space of rank~$r$ to~$V$; for the spaces of rank \m{r+2\nu},
\m{\nu\ge1}, the corresponding absolute value is larger by the factor~\m{q_{w}^{\nu}}. So the theorem says that a
function with this Hecke root and without descent to~$W$ has no descent to the higher spaces of the tower
either, in the range \m{2r+2\nu+3<n}. In the application (§\ref{s:use}) Theorems 1 and~2 are applied one after
the other, to \m{Y=\conj{\coef{L}{\Xi}}}.

\sect{The setting}{s:setting}
\begin{conventions}\phantomsection\label{c:conventions}
The conventions of the second note of this sequence hold: $E$ is a CM field with a CM type; $\tau$ denotes its
infinite places and $w$ its finite places, with conjugate~\m{\conj{w}}, completion~$E_{w}$, integers~$\oo_{w}$,
residue field~$\FF_{w}$ with $q_{w}$ elements, a fixed uniformiser~$\varpi_{w}$ and the absolute value
\m{|\cdot|_{w}} with \m{|\varpi_{w}|_{w}<1}; \m{\AAf} is the ring of finite adeles and \m{\ohat} the product of
the~$\oo_{w}$. The index sets $[r]$, $[n]$, $[r']$, $[\nu]$ are finite sets with $r$, $n$, $r'$, $\nu$ elements;
an \m{n\times r} matrix has rows indexed by~$[n]$ and columns by~$[r]$. For a matrix~$X$, \m{X^{*}} is the
conjugate transpose and \m{X^{\mathsf{T}}} the transpose. A complex vector space is written~$\mathcal{V}$ when it
carries a representation. Where a statement says so, $E$ carries Weil data (Setting~\ref{s:weil}). Letters of
the handout's overview that are not displayed with the data as an argument (among them $\Upsilon$, $f$,
$\varepsilon$, $\beta$, $\delta$, $n$ and~$V$) have, in this note, only the meaning given in its notation list.
\end{conventions}

\begin{setting}[groups, the Weil representation, descents; quoted, as in the second note]\label{s:weil}
The field~$E$ carries Weil data; of these we use the splitting character and its restriction~$\upsilon_{1}$ to
the ideles of norm one. For a family of signs~$\varepsilon$ on~$[r]$, \m{\Gad{\varepsilon}=\Ginf{\varepsilon}\times\Gfin{r}}
is the adelic unitary group: \m{\Ginf{\varepsilon}} is the product over the infinite places of the unitary groups
of \m{\operatorname{diag}(\varepsilon)}, and \m{\Gfin{r}} the group of \m{u\in\mathrm{GL}_{r}(\AAf)} with \m{u^{*}u=1}.
The group \m{\Kz{r}=\Gfin{r}\cap\mathrm{GL}_{r}(\ohat)} is open and compact, and for \m{w\ne\conj{w}},
\m{\red{w}\colon\Kz{r}\to\mathrm{GL}_{r}(\FF_{w})} is the reduction. For a set~$S$ of finite places, an element of
\m{\Gfin{r}} is \emph{supported on~$S$} if its components at all places outside~$S$ and not conjugate to a place
of~$S$ are~$1$. A hermitian space with a point, of signs~$\varepsilon$, has a matrix~$J$ over~$E$, a group of
rational points embedded in \m{\Gad{\varepsilon}}, and a lattice at every point of \m{\Gad{\varepsilon}}; the
standard definite space \m{\std{E}{n}} has signs~\m{\one{n}}. A function on \m{\Gad{\varepsilon}} is \emph{automorphic}
on such a space if it is continuous, left invariant under the rational points and right invariant under an
open subgroup of \m{\Gfin{r}}. An automorphic character is a continuous unitary character of the adeles of norm
one of~$E$, trivial on the rational ones; its \emph{base change} is the Hecke character \m{z\mapsto\chi(z/\conj{z})}
of~$E$, and a Hecke character is unramified at~$w$ if it is trivial on the units of~$\oo_{w}$.

For spaces $W$,~$V$ of signs $\varepsilon$ on~$[r]$ and $\varepsilon'$ on~$[n]$ and automorphic characters
$\chi$,~$\chi'$, the library provides the Banach space \m{\mathcal{S}(W,V)} of data on \m{\operatorname{Hom}(W,V)},
with the Heisenberg operators \m{\varrho(z)} of adelic vectors~$z$ of \m{\operatorname{Hom}(W,V)}, the Weil
representation \m{\omega^{W,V}_{\chi,\chi'}} of \m{\Gad{\varepsilon}\times\Gad{\varepsilon'}} on it, and the theta
functional \m{\theta_{W,V}}, invariant under the rational points of both groups. The vacuum
\m{\vacf\in\mathcal{S}(W,V)} is the Gaussian at the infinite places times the characteristic function of the
integral matrices at the finite places. For a Fock vector~$v$ of the pair and a datum~$f$ there is the datum
\m{v\otimes f}; \m{v^{\varepsilon,\varepsilon'}_{p}} is the Fock vector of a polynomial~$p$ in the variables
\m{X_{\tau,i,c}}. The theta descent of a function~$Y$ on \m{\Gad{\varepsilon'}} with datum~$f$ is the function
\[
\Desc{W}{V}{\chi}{\chi'}(f,Y)(u)=\int_{[V]}\theta_{W,V}\bigl(\omega^{W,V}_{\chi,\chi'}(u,h)f\bigr)\,Y(h)\,\dd h
\qquad(u\in\Gad{\varepsilon}),
\]
the integral over the quotient of \m{\Gad{\varepsilon'}} by the rational points of~$V$, with the measure of the
second note. In this note the second character is always the trivial character~$1$.
\end{setting}
\leanline{in the background library: \lnm{WeilData}, \lnm{WeilData.upsilonU1}, \lnm{G}, \lnm{Gf}, \lnm{Kf},
\lnm{Gf.supportedOn}, \lnm{Space}, \lnm{stdSpace}, \lnm{Space.IsAutomorphic}, \lnm{AutChar},
\lnm{AutChar.baseChange}, \lnm{CharUnramifiedAt}, \lnm{pairWeil}, \lnm{pairThetaKer}, \lnm{Space.lift},
\lnm{Space.integrated}}

\begin{definition}[data unramified at a set of places; vanishing descents; quoted]\label{d:dv}
Let $W$,~$V$ be as in Setting~\ref{s:weil} and $S$ a set of finite places of~$E$. The operator \m{\PS{S}} on
\m{\mathcal{S}(W,V)} is the mean of the Heisenberg operators \m{\varrho(z)} over the integral finite vectors~$z$
whose components vanish at the places that are neither in~$S$ nor conjugate to a place of~$S$; it is a
projector, and a datum~$f$ with \m{\PS{S}f=f} is called \emph{unramified at~$S$}: at the places of~$S$ and their
conjugates it behaves as the vacuum. For an automorphic character~$\chi$ and a function~$Y$ on
\m{\Gad{\varepsilon'}}, the statement \m{\DV{W}{V}{\chi}{S}{Y}} is:
\[
\Desc{W}{V}{\chi}{1}(f,Y)=0\qquad\text{for every } f\in\mathcal{S}(W,V)\text{ with }\PS{S}f=f .
\]
\end{definition}
\leanline{in the background library: \lnm{Space.projOn}, \lnm{DescentsVanish}}

\begin{setting}[auxiliary pairs; Input from the unit on the rows]\label{s:aux}
Let $E$ carry Weil data. An auxiliary pair~$O$ over~$E$ for the index sets $[r]$ and~$[n]$, \m{r<n}, is as in
the second note. We use: its level system~$\SigO{O}$ with the chosen places~$\MO{O}$ and the twist places~$\QO{O}$,
and \m{\plO{O}=\MO{O}\cup\QO{O}}, a finite set of places none of which is conjugate to another or to itself; its
space \m{\WO{O}} of signs~$\sgO{O}$ on~$[r]$; the injection \m{\incO{O}\colon[r]\to[n]} and the \m{n\times r}
matrix \m{\IncM{O}} with entries~$1$ at the places \m{(\incO{O}(c),c)} and $0$ elsewhere; the multiplicative
characters \m{\psiO{O}{w}} of the fields~$\FF_{w}$, extended by~$0$, trivial for \m{w\notin\plO{O}}; the weight
\m{(\kap{O},\eO{O})}, the space \m{\Vk{n}{\kap{O}}} and the representation \m{\rhoW{\kap{O}}{\eO{O}}} of \m{\Ginf{\one{n}}} on
it; the character \m{\chO{O}}, a power of~$\upsilon_{1}$; the Fock vectors \m{\pvO{O}{p}=v^{\sgO{O},\one{n}}_{p}}; and
the following.
\begin{enumerate}[label=\textup{(\alph*)},leftmargin=2.2em,itemsep=2pt,topsep=2pt]
\item \emph{Level.} \m{\LO{O}} is the group of families \m{(j_{w})_{w\in\plO{O}}}, \m{j_{w}\in\mathrm{GL}_{n}(\FF_{w})},
with \m{j_{w}\IncM{O}=\IncM{O}} for \m{w\in\MO{O}} and, for \m{w\in\QO{O}}, with vanishing entries in the rows of the
image of~$\incO{O}$ and the columns outside it; \m{\etaO{O}(j)} is the product over \m{w\in\QO{O}} of
\m{\psiO{O}{w}(\det(\incO{O}\text{-block of }j_{w}))^{-1}}; \m{\KO{O}} is the group of \m{g'\in\Kz{n}} with
\m{(\red{w}g')_{w}\in\LO{O}}, and \m{\tetaO{O}(g')=\etaO{O}((\red{w}g')_{w})}.
\item \emph{Forms.} An \emph{algebraic form of weight} \m{(\kap{O},\eO{O})} is a function
\m{\Xi\colon\Gfin{n}\to\Vk{n}{\kap{O}}} with \m{\Xi(\gamma_{\mathrm{f}}h)=\rhoW{\kap{O}}{\eO{O}}(\gamma)\Xi(h)} for rational
points~$\gamma$ of \m{\std{E}{n}}, whose stabiliser under right translation is open; \m{\Gfin{n}} acts on these
forms by right translation, written \m{\mathrm{R}(g')\Xi=\Xi(\,\cdot\,g')}. The space \m{\FormO{O}} consists of the
algebraic forms with \m{\Xi(hg')=\tetaO{O}(g')\Xi(h)} for \m{g'\in\KO{O}}. For a linear form~$L$ on \m{\Vk{n}{\kap{O}}},
\m{\coef{L}{\Xi}} is the function \m{h\mapsto L\bigl(\rhoW{\kap{O}}{\eO{O}}(h_{\infty}^{-1})\Xi(h_{\mathrm{f}})\bigr)} on
\m{\Gad{\one{n}}}.
\item \emph{The action of the chosen places.} Let $\pi$ be a representation of \m{\Gfin{n}} on a complex vector
space~$\mathcal{V}$, and let \m{\mathcal{V}^{O}} be the subspace of the~$\phi$ with
\m{\pi(g')\phi=\tetaO{O}(g')\phi} for all \m{g'\in\KO{O}}. The group
\m{\ChGL{\SigO{O}}{r}=\prod_{w\in\MO{O}}\mathrm{GL}_{r}(\FF_{w})} acts on \m{\mathcal{V}^{O}}: a family \m{A=(A_{w})} acts
by
\[
\chact{O}{\pi}{A}\,\phi=\prod_{w\in\MO{O}}\psiO{O}{w}(\det A_{w})\cdot\pi(g_{A})\,\phi .
\]
Here \m{g_{A}\in\Kz{n}} is the element that the library chooses with
\[
\red{w}g_{A}=\operatorname{diag}(A_{w},1)\quad(w\in\MO{O}),\qquad\red{w}g_{A}=1\quad(w\in\QO{O}),
\]
and \m{\operatorname{diag}(A_{w},1)} acts by~$A_{w}$ on the
basis vectors with indices in the image of~$\incO{O}$ and fixes the others. For \m{\mathcal{V}} the space of
algebraic forms of~(b) with \m{\pi=\mathrm{R}}, we have \m{\mathcal{V}^{O}=\FormO{O}}, and
\m{\chact{O}{\mathrm{R}}{A}} is the action on forms.
\item \emph{The finite datum.} \m{\PsiO{O}{X}}, for a family \m{X=(X_{w})_{w\in\plO{O}}} of \m{n\times r} matrices
over the fields~$\FF_{w}$, is \m{\prod_{w\in\QO{O}}\psiO{O}{w}(\det(\incO{O}\text{-rows of }X_{w}))} if \m{X_{w}=\IncM{O}}
for all \m{w\in\MO{O}}, and $0$ otherwise. The finite datum \m{\Ups{O}\in\mathcal{S}(\WO{O},\std{E}{n})} has, at the
places of \m{\plO{O}}, the component \m{X\mapsto\PsiO{O}{(X\bmod w)_{w}}} on the integral matrices and $0$ off them, and
is the vacuum at all other places.
\end{enumerate}
\end{setting}
\leanline{in the unit on the rows: \lnm{AuxPair}, \lnm{auxLevel}, \lnm{AuxPair.auxLevelChar},
\lnm{AuxPair.levelGroup}, \lnm{AuxPair.levelChar}, \lnm{AuxPair.FormV}, \lnm{AuxPair.chosenGLRep},
\lnm{AuxPair.formAction}, \lnm{AuxPair.resDetDatum}, \lnm{AuxPair.auxDatum}, \lnm{AuxPair.a}\sep in the
background library: \lnm{algFormCoeff}, \lnm{levelRep}, \lnm{liftResidue}}

\begin{definition}[reflection sums; quoted]\label{d:refl}
Let $\FF$ be a finite field, $\psi$ a multiplicative character of~$\FF$ extended by~$0$, and \m{y\in\FF^{r}}. The
\emph{reflection sum} is the element of the complex group algebra of \m{\mathrm{GL}_{r}(\FF)}
\[
\Refl{\psi}{y}=\sum_{A}\psi(\det A)^{-1}\,A ,
\]
the sum over the invertible matrices of the form \m{A=1+\xi\,y^{\mathsf{T}}} with a column vector
\m{\xi\in\FF^{r}}. For finite groups $G_{w}$ and elements \m{a_{w}} of their group algebras, \m{\bigotimes_{w}a_{w}} is
the element of the group algebra of \m{\prod_{w}G_{w}} whose coefficient at \m{(A_{w})_{w}} is the product of the
coefficients of the~$a_{w}$ at the~$A_{w}$.
\end{definition}
\leanline{in the background library: \lnm{reflectionSum}, \lnm{piTensorElem}}

For \m{y\ne0} the invertible matrices \m{1+\xi y^{\mathsf{T}}} are the elements of \m{\mathrm{GL}_{r}(\FF)} that fix
every vector of the hyperplane on which the linear form \m{y^{\mathsf{T}}} vanishes: the transvections and the
pseudo-reflections along it;
for \m{y=0} the sum is the identity. Hypothesis~(b) of Theorem~1 is a property that the consumers of the
theorem establish for the forms to which they apply it.

\begin{definition}[Hecke polynomials at a split place; quoted]\label{d:hecke}
Let $w$ be a finite place of~$E$ with \m{w\ne\conj{w}}, and $\pi$ a representation of \m{\Gfin{r'}} on a complex
vector space~$\mathcal{V}$. Since $w$ is split, \m{\mathrm{GL}_{r'}(E_{w})} embeds in \m{\Gfin{r'}}: a matrix~$g$ goes
to the element with component~$g$ at~$w$, \m{(g^{*})^{-1}} at~\m{\conj{w}}, and $1$ elsewhere. A vector is
\emph{spherical at~$w$} if it is fixed by the elements of \m{\Kz{r'}} supported on~\m{\{w\}}. On such vectors the
Hecke operators \m{T_{w,a}}, \m{0\le a\le r'}, are those of the double cosets of
\m{\operatorname{diag}(\varpi_{w}1_{a},1_{r'-a})} in \m{\mathrm{GL}_{r'}(E_{w})}, and the \emph{Hecke polynomial} at a
complex number~$T$ is the operator
\[
\Hk{\pi}{w}{T}=\sum_{a=0}^{r'}(-T)^{a}\,q_{w}^{\,a(a-1)/2}\;T_{w,a}.
\]
For signs~$\varepsilon'$ on~$[r']$, \m{\Rt{\varepsilon'}} is the representation of \m{\Gfin{r'}} on the complex
functions on \m{\Gad{\varepsilon'}} by right translation.
\end{definition}
\leanline{in the background library: \lnm{FinPlace.IsSplit}, \lnm{atSplit}, \lnm{sphericalAt},
\lnm{heckeSplit}, \lnm{heckePolyOp}, \lnm{finRightRep}}

Up to the normalisation of the parameters, \m{\Hk{\pi}{w}{T}} acts on the spherical vector of an unramified
principal series of \m{\mathrm{GL}_{r'}(E_{w})} by the product of the factors \m{1-\alpha_{i}T} over its Satake
parameters~$\alpha_{i}$: it is the reciprocal of the standard local Euler factor in the variable~$T$. For the
trivial representation the factors are \m{1-q_{w}^{a}T}, \m{0\le a<r'}, which is the second part of (Q7) below.

\begin{definition}[splittings and zero terms; quoted]\label{d:split}
Let $W$ and~$W'$ be hermitian spaces with points, of signs $\varepsilon$ on~$[r]$ and $\varepsilon'$ on~$[r']$,
with matrices $J$ and~$J'$. A \emph{splitting}~$\Pi$ of~$W'$ over~$W$ with planes indexed by~$[\nu]$ consists of
matrices \m{k_{\Pi}}, \m{l_{\Pi}} of size \m{r'\times\nu} and \m{i_{\Pi}} of size \m{r'\times r} over~$E$ with
\m{k_{\Pi}^{*}J'k_{\Pi}=0}, \m{l_{\Pi}^{*}J'l_{\Pi}=0}, \m{l_{\Pi}^{*}J'k_{\Pi}=1}, \m{i_{\Pi}^{*}J'i_{\Pi}=J},
\m{i_{\Pi}^{*}J'k_{\Pi}=0}, \m{i_{\Pi}^{*}J'l_{\Pi}=0}, and it is required that \m{r'=r+2\nu}. So $W'$ is the
orthogonal sum of~$W$ and of $\nu$ hyperbolic planes with isotropic bases given by the columns of \m{k_{\Pi}}
and~\m{l_{\Pi}}; we write \m{W'\cong\hyp{W}{\nu}}. The splitting is \emph{adapted} if, in the coordinates of the
points of $W$ and~$W'$, the frame \m{(l_{\Pi},k_{\Pi},i_{\Pi})} has orthonormal columns at the infinite places
and integral entries at the finite places. For a space~$V$ of signs~$\one{n}$, the \emph{zero term} of a datum
along~$\Pi$ is given by the bounded linear map
\[
\ZT{\Pi}{V}{0}\colon\mathcal{S}(W',V)\to\mathcal{S}(W,V),
\]
which in a Schrödinger model takes the mean of a function over the adelic translations along the maps that
vanish on the orthogonal complement of the columns of~\m{k_{\Pi}}, modulo the rational ones, and restricts to
\m{\operatorname{Hom}(W,V)}.
\end{definition}
\leanline{in the background library: \lnm{Space.Splitting}, \lnm{IsBaseAdapted},
\lnm{Space.Splitting.homSource}, \lnm{Space.Splitting.zeroTerm}, \lnm{Space.Splitting.sourceTerm} (at the
argument~0)}

\statement{inksoft}{\textbf{Quoted from the background library, I: data and descents.} These facts are taken for
granted; references and their state are in §\ref{s:notes}, note~6. $W$, $V$, $\varepsilon$, $\varepsilon'$, $\chi$
are as in Setting~\ref{s:weil}, with \m{\varepsilon'=\one{n}}, and $\omega$ stands for \m{\omega^{W,V}_{\chi,1}}.

\vspace{3pt}\textbf{(Q1)} \emph{The descent as a functional.} Let $Y$ be automorphic on~$V$. Then $Y$ is bounded,
the quotient~$[V]$ has finite measure, and for every \m{u\in\Gad{\varepsilon}} the map \m{f\mapsto\Desc{W}{V}{\chi}{1}(f,Y)(u)}
is a continuous linear functional on \m{\mathcal{S}(W,V)}. For \m{u'\in\Gad{\varepsilon}} and \m{h'\in\Gad{\one{n}}},
\m{\Desc{W}{V}{\chi}{1}(f,Y)(uu')=\Desc{W}{V}{\chi}{1}\bigl(\omega(u',h')f,\,Y(\,\cdot\,h')\bigr)(u)}, and
\m{\Desc{W}{V}{\chi}{1}(f,Y)(\gamma u)=\Desc{W}{V}{\chi}{1}(f,Y)(u)} for rational points~$\gamma$ of~$W$.

\vspace{3pt}\textbf{(Q2)} \emph{Density.} Let $S$ be a set of finite places and $f$ a datum with \m{\PS{S}f=f}. If
a continuous linear functional on \m{\mathcal{S}(W,V)} vanishes on all data
\m{v^{\varepsilon,\one{n}}_{p}\otimes\varrho(z)\vacf}, with $p$ a polynomial and $z$ a finite vector whose components at
the places of~$S$ and at their conjugates are zero, then it vanishes on~$f$. These data are unramified at~$S$.

\vspace{3pt}\textbf{(Q3)} \emph{Functions at a finite set of places.} Let $P$ be a finite set of finite places,
none conjugate to another or to itself, and \m{E_{P}=\prod_{w\in P}E_{w}}. A Schwartz--Bruhat function~$\Psi$ on the
\m{n\times r} matrices over~$E_{P}$ defines a datum \m{\Psi\cdot\vacf} that has the component~$\Psi$ at~$P$ and is the
vacuum elsewhere; it is linear in~$\Psi$. Let \m{\operatorname{Box}} be the set of matrices with entries in
\m{\prod_{w\in P}\oo_{w}}; for an idele~$\mathfrak{t}$ with integral components and a matrix~$z$,
\m{\cosI{\mathfrak{t}}{z}} is the characteristic function of \m{z+\mathfrak{t}\operatorname{Box}}. (a) For a finite vector~$z$
supported on~$P$ and its conjugates, \m{\varrho(z)\vacf} is a complex multiple of \m{\Psi\cdot\vacf} for a
Schwartz--Bruhat function~$\Psi$ at~$P$. (b) A linear space of functions that contains all
\m{\cosI{\mathfrak{t}}{z}} contains every Schwartz--Bruhat function at~$P$. (c) If such a function~$\Psi$ is
invariant under the translations by \m{\mathfrak{t}\operatorname{Box}}, then
\m{\Psi=\sum_{z'\in\mathcal{Z}}\Psi(z')\,\cosI{\mathfrak{t}}{z'}} for a finite set~$\mathcal{Z}$ of matrices with
\m{\Psi(z')\ne0}. (d) A function of the residues of an integral matrix at the places of~$P$, extended by~$0$, is
Schwartz--Bruhat, and takes the same value at two matrices whose difference has entries of absolute value
\m{<1} at all \m{w\in P}.

\vspace{3pt}\textbf{(Q4)} \emph{Finite elements acting on such data.} Let $\chi$ be a power of~$\upsilon_{1}$,
$v$ a Fock vector and $\Psi$ as in (Q3). (a) For \m{g'\in\Kz{n}} supported on~$P$:
\m{\omega(1,(1,g'))(v\otimes\Psi\cdot\vacf)=v\otimes\Psi'\cdot\vacf} with \m{\Psi'(X)=\Psi(g'^{-1}X)}. (b) For
\m{\alpha\in\Gfin{r}} supported on~$P$: \m{\omega((1,\alpha),1)(v\otimes\Psi\cdot\vacf)} is a non-zero complex multiple
of \m{v\otimes\Psi''\cdot\vacf} with \m{\Psi''(X)=\Psi(X\alpha)}.

\vspace{3pt}\textbf{(Q5)} \emph{Scalar functions of forms.} For an auxiliary pair~$O$, an algebraic form~$\Xi$ of
weight \m{(\kap{O},\eO{O})} and a linear form~$L$ on \m{\Vk{n}{\kap{O}}}: \m{\coef{L}{\Xi}} is automorphic on
\m{\std{E}{n}}, and \m{\coef{L}{(\mathrm{R}(g')\Xi)}(h)=\coef{L}{\Xi}(hg')} for \m{g'\in\Gfin{n}}.

\vspace{3pt}\textbf{(Q6)} \emph{Primitive matrices.} Let $w$ be a finite place. An \m{n\times r} matrix over~$E_{w}$
with a left inverse becomes, after multiplication on the right by a suitable element of
\m{\mathrm{GL}_{r}(E_{w})}, a matrix with entries in~$\oo_{w}$ that has a left inverse with entries in~$\oo_{w}$. A
square matrix over~$\oo_{w}$ whose reduction is invertible is invertible over~$\oo_{w}$.}
\leanline{in the background library: (Q1) \lnm{ThetaKernel.exists\_clm\_lift}, \lnm{ThetaKernel.lift\_equivariant},
\lnm{lift\_mul\_right}, \lnm{ThetaKernel.fixesLeft\_lift}\sep (Q2) \lnm{Space.ext\_integrated}\sep (Q3)
\lnm{Space.name}, \lnm{cosetIndicator}, \lnm{Space.name\_schrodingerAct\_boxIndicator},
\lnm{IsSchwartzBruhat.mem\_of\_cosets}, \lnm{IsSchwartzBruhat.eq\_sum\_cosetIndicator},
\lnm{isSchwartzBruhat\_residueDatum}, \lnm{residueDatum\_congr}\sep (Q4) \lnm{pairWeil\_name\_of\_mem},
\lnm{pairWeil\_eq\_smul\_of\_projOn\_eq}, \lnm{pairWeil\_name\_on}, \lnm{pairWeil\_integrated\_fin}\sep (Q5)
\lnm{Space.isAutomorphic\_algFormCoeff}, \lnm{Space.algFormCoeff\_algFormRep}\sep (Q6)
\lnm{IsValuationPair.exists\_isPrimitiveMatrix\_mul}, \lnm{IsPrimitiveMatrix.exists\_map\_subtype\_eq},
\lnm{isUnit\_of\_mapMatrix}}

\statement{inksoft}{\textbf{Quoted from the background library, II: Hecke operators.}
Here $W'$ has signs~$\varepsilon'$ on~$[r']$, $V$ has signs~$\one{n}$, $\chi_{1}$ is an automorphic character and
$\omega'$ stands for \m{\omega^{W',V}_{\chi_{1},1}}.

\vspace{3pt}\textbf{(Q7)} \emph{Hecke polynomials as sums over subspaces.} Let \m{w\ne\conj{w}}. There are elements
\m{g_{U}\in\Gfin{r'}}, one for each linear subspace~$U$ of \m{\FF_{w}^{r'}}, supported on~\m{\{w\}}, such that for every
representation~$\pi$ of \m{\Gfin{r'}}, every vector~$\phi$ spherical at~$w$ and every complex~$T$:
\m{\Hk{\pi}{w}{T}\phi=\sum_{U}(-T)^{\dim U}q_{w}^{\dim U(\dim U-1)/2}\,\pi(g_{U})\phi}. Moreover (the $q$-binomial
theorem) \m{\sum_{U}(-T)^{\dim U}q_{w}^{\dim U(\dim U-1)/2}=\prod_{0\le a<r'}(1-q_{w}^{a}T)}.

\vspace{3pt}\textbf{(Q8)} \emph{Adjoints.} Let \m{w\ne\conj{w}}, let $\pi$,~$\pi'$ be representations of
\m{\Gfin{r'}} on $\mathcal{V}$,~$\mathcal{V}'$ and $\beta$ a bilinear map on \m{\mathcal{V}\times\mathcal{V}'} with
\m{\beta(\pi(g)\phi,\pi'(g)\phi')=\beta(\phi,\phi')}. For vectors spherical at~$w$:
\m{\beta(\Hk{\pi}{w}{T}\phi,\phi')=\beta(\phi,\Hk{\pi'}{\conj{w}}{T}\phi')}. A linear map that commutes with the elements
supported on~\m{\{w\}} commutes with the Hecke polynomials at~$w$. Vectors spherical at~$w$ are spherical
at~\m{\conj{w}}.

\vspace{3pt}\textbf{(Q9)} \emph{Hecke transfer.} Let \m{w\ne\conj{w}} and \m{r'\le n}. Suppose that the base change
of~$\chi_{1}$ is unramified at~$w$ and that $f$ is unramified at~\m{\{w\}}. Let $\Lambda$ be a linear map from
\m{\mathcal{S}(W',V)} to a complex vector space with commuting representations $\pi$ of \m{\Gfin{r'}} and $\rho$ of
\m{\Gfin{n}}, such that \m{\Lambda(\omega'((1,g),(1,h))f)=\pi(g)\rho(h)\Lambda f} for all \m{g\in\Gfin{r'}},
\m{h\in\Gfin{n}}. Then \m{\Lambda f} is spherical at~$w$ for $\pi$ and for~$\rho$, and for every complex~$T$
\[
\Hk{\rho}{w}{T}\,\Lambda f=\prod_{0\le a<n-r'}\bigl(1-q_{w}^{a}\,A\,T\bigr)\cdot\Hk{\pi}{\conj{w}}{q_{w}^{\,n-r'}AB^{-1}T}\,\Lambda f ,
\]
with complex numbers $A$,~$B$ that depend on $\chi_{1}$, $w$, $r'$, $n$ and have \m{|A|=q_{w}^{r'/2}},
\m{|B|=q_{w}^{n/2}}. The character~$\upsilon_{1}$ has unramified base change at every split place.}

\statement{inksoft}{\textbf{Quoted from the background library, III: zero terms, growth.} The conventions are
those of~II.

\vspace{3pt}\textbf{(Q10)} \emph{Zero terms.} Let $\Pi$ be an adapted splitting of~$W'$ over~$W$ with planes indexed
by~$[\nu]$, and \m{\chi_{1}=\chi\upsilon_{1}^{\nu}}. (a) For every \m{u\in\Gad{\varepsilon'}} there are an element~$g$ of
the standard maximal compact subgroup of \m{\Gad{\varepsilon'}}, with finite component in \m{\Kz{r'}}, an element
\m{u_{0}\in\Gad{\varepsilon}} and a complex number~$a$ such that for all \m{h\in\Gad{\one{n}}} and all data~$f$:
\m{\ZT{\Pi}{V}{0}\bigl(\omega'(u,h)f\bigr)=a\cdot\omega^{W,V}_{\chi,1}(u_{0},h)\,\ZT{\Pi}{V}{0}\bigl(\omega'(g^{-1},1)f\bigr)}. (b) If
\m{\PS{S}f=f}, then \m{\PS{S}f_{1}=f_{1}} for \m{f_{1}=\ZT{\Pi}{V}{0}(\omega'(g^{-1},1)f)}. (c) If \m{\nu=0}, then
\m{\theta_{W,V}\circ\ZT{\Pi}{V}{0}=\theta_{W',V}}. (d) If \m{\nu\ge1}, there are a space~$W_{1}$ with an adapted
splitting over~$W$ with \m{\nu-1} planes, and an adapted splitting~$\Pi'$ of~$W'$ over~$W_{1}$ with one plane.

\vspace{3pt}\textbf{(Q11)} \emph{Growth and decay.} Every point~$u$ of an adelic unitary group has a height
\m{\operatorname{ht}(u)\ge1}, invariant under rational points on the left; on a definite space it is~$1$. (a) There
are constants $C$,~$a$ with \m{|\theta_{W',V}(\omega'(u,h)\varphi)|\le C\lVert\varphi\rVert\operatorname{ht}(u)^{a}} for all
$u$,~$h$,~$\varphi$. (b) If \m{\operatorname{ht}(u)>2}, the lattice of~$W'$ at~$u$ contains a non-zero isotropic
vector of length at most \m{2/\operatorname{ht}(u)}; rational frames with the same Gram matrix are conjugate
under the rational points (Witt); and the lattice of~$V$ at every~$h$ has a basis of bounded length. (c) Let
$\Pi'$ be an adapted splitting of~$W'$ over~$W_{1}$ with one plane and
\m{\thd=\theta_{W',V}-\theta_{W_{1},V}\circ\ZT{\Pi'}{V}{0}}. Then $\thd$ is invariant under the Heisenberg operators
of the adelic translations along the isotropic vector~\m{k_{\Pi'}}, and \m{|\thd(\varphi)|} is at most the supremum
of \m{|\theta_{W',V}(\varphi)-\theta_{W',V}(\varrho(z)\varphi)|} over these translations~$z$. (d) \emph{The decay
estimate, in the case used.} For data \m{f=v^{\varepsilon',\one{n}}_{p}\otimes\varrho(z)\vacf} there is a constant~$C_{0}$
with \m{|\thd(\omega'(u,h)f)|\le C_{0}} for all $h$ and all~$u$ such that either \m{\operatorname{ht}(u)\le2}, or
the isotropic vector~\m{k_{\Pi'}} is, in the lattice of~$W'$ at~$u$, of length at most \m{2/\operatorname{ht}(u)}.}
\leanline{in the background library: (Q7) \lnm{heckePolyOp\_eq\_sum}, \lnm{q\_binomial\_submodule}\sep (Q8)
\lnm{heckePolyOp\_adjoint}, \lnm{map\_heckePolyOp}, \lnm{Gf.supportedOn\_star}\sep (Q9)
\lnm{hecke\_transfer\_of\_equivariant}, \lnm{WeilData.charUnramifiedAt\_upsilon}\sep (Q10)
\lnm{Space.Splitting.iwasawa}, \lnm{Space.Splitting.sourceTerm\_radicalElt},
\lnm{Space.Splitting.sourceTerm\_levi}, \lnm{projOn\_pairWeil},
\lnm{Space.Splitting.IsAdapted.omega\_projOn\_of\_isMixedLaw}, \lnm{exists\_isBaseAdapted\_stage},
\lnm{removePlanes}\sep (Q11) \lnm{ThetaKernel.moderate\_growth}, \lnm{Space.exists\_short\_isotropic},
\lnm{Space.witt}, \lnm{Space.exists\_short\_basis\_matrix}, \lnm{Space.Splitting.zeroTerm\_rho\_kAdelic},
\lnm{norm\_quotMean\_le}, \lnm{exists\_norm\_name\_le}}

\sect{Three lemmas on the level}{s:level}
\begin{lemma}[prescribed components at the places of the pair]\label{l:atplace}
Let $O$ be an auxiliary pair (Setting~\ref{s:aux}), $[r']$ a finite index set, and \m{g=(g_{w})_{w\in\plO{O}}} a
family of matrices \m{g_{w}\in\mathrm{GL}_{r'}(E_{w})}. Then there is an element of \m{\Gfin{r'}}, supported on
\m{\plO{O}}, whose component at every \m{w\in\plO{O}} is~$g_{w}$.
\end{lemma}
\begin{proof}
For \m{w\in\plO{O}} we have \m{w\ne\conj{w}}, so \m{\mathrm{GL}_{r'}(E_{w})} embeds in \m{\Gfin{r'}}
(Definition~\ref{d:hecke}): the image of~$g_{w}$ has the component~$g_{w}$ at~$w$ and is supported on~\m{\{w\}}. No place
of \m{\plO{O}} is conjugate to another, so the image of~$g_{w}$ has the component~$1$ at the other places of
\m{\plO{O}}, and the product of these images over \m{w\in\plO{O}} is the element sought.
\end{proof}
\leanline{\lean{existsatPlaceeq}}

\begin{lemma}[an annihilator, or a congruence to the inclusion matrix]\label{l:annihilator}
Let \m{\iota\colon[r]\to[n]} be an injective map, $w$ a finite place of~$E$, and $X$ an \m{n\times r} matrix
over~$E_{w}$ that has a left inverse over~$E_{w}$. Let \m{\operatorname{Inc}_{\iota}} be the \m{n\times r} matrix with
entries~$1$ at the places \m{(\iota(c),c)} and $0$ elsewhere. Then there is an \m{n\times n} matrix~$\ell$ with
entries in~$\oo_{w}$ such that \m{\ell X=0} and such that the following holds: if all entries
\m{\ell_{i,\iota(c)}}, \m{i\in[n]}, \m{c\in[r]}, have absolute value \m{<1}, then there is
\m{\alpha\in\mathrm{GL}_{r}(E_{w})} such that every entry of \m{X\alpha-\operatorname{Inc}_{\iota}} has absolute value
\m{<1}.
\end{lemma}
\begin{proof}
By (Q6) there is \m{\beta\in\mathrm{GL}_{r}(E_{w})} such that \m{Z=X\beta} has entries in~$\oo_{w}$ and a left
inverse~$Z'$ with entries in~$\oo_{w}$: \m{Z'Z=1}. Put \m{\ell=ZZ'-1}. Then
\m{\ell X=(ZZ'-1)Z\beta^{-1}=0}. Suppose that all \m{\ell_{i,\iota(c)}} have absolute value \m{<1}, and let $\beta'$ be
the \m{r\times r} matrix with the entries \m{Z'_{c',\iota(c)}}. The entry of \m{Z\beta'} at \m{(i,c)} is
\m{(ZZ')_{i,\iota(c)}}, which is congruent modulo the maximal ideal to the entry of the identity matrix at
\m{(i,\iota(c))}, that is, to the entry of \m{\operatorname{Inc}_{\iota}} at \m{(i,c)}. In particular the product of
the \m{\iota}-rows of~$Z$ with~$\beta'$ reduces to the identity, so $\beta'$ is invertible modulo the maximal
ideal, hence over~$\oo_{w}$ by (Q6). Then \m{\alpha=\beta\beta'} has \m{X\alpha=Z\beta'}, with the congruence just
shown.
\end{proof}
\leanline{\lean{existsannihilatororcongrinclMatrix}}

\begin{lemma}[an integral element with block reductions acts through the chosen places]\label{l:level}
Let $O$ be an auxiliary pair (Setting~\ref{s:aux}), $\pi$ a representation of \m{\Gfin{n}} on a complex vector
space~$\mathcal{V}$, and \m{\phi\in\mathcal{V}^{O}}. Let \m{g\in\Kz{n}}, with reductions \m{\conj{g}_{w}=\red{w}g} for
\m{w\in\plO{O}}, and let \m{s=(s_{w})_{w\in\plO{O}}} be a family of matrices \m{s_{w}\in\mathrm{GL}_{r}(\FF_{w})}. Suppose
that \m{\conj{g}_{w}\,\IncM{O}=\IncM{O}\,s_{w}} for every \m{w\in\MO{O}} and
\m{\conj{g}_{w}^{\mathsf{T}}\,\IncM{O}=\IncM{O}\,s_{w}^{\mathsf{T}}} for every \m{w\in\QO{O}}. Then
\[
\pi(g)\,\phi=\Bigl(\prod_{w\in\plO{O}}\psiO{O}{w}(\det s_{w})\Bigr)^{-1}\chact{O}{\pi}{(s_{w})_{w\in\MO{O}}}\,\phi .
\]
\end{lemma}
\begin{proof}
Let \m{A=(s_{w})_{w\in\MO{O}}} and let \m{g_{A}\in\Kz{n}} be as in Setting~\ref{s:aux}(c). The reductions of
\m{g_{A}^{-1}g} are \m{j_{w}=\operatorname{diag}(s_{w},1)^{-1}\conj{g}_{w}} at \m{w\in\MO{O}} and \m{j_{w}=\conj{g}_{w}} at
\m{w\in\QO{O}}. At \m{w\in\MO{O}}:
\[
j_{w}\,\IncM{O}=\operatorname{diag}(s_{w},1)^{-1}\,\IncM{O}\,s_{w}=\IncM{O}\,s_{w}^{-1}s_{w}=\IncM{O}.
\]
At \m{w\in\QO{O}}: the hypothesis says that the transpose of the rows of~$\conj{g}_{w}$ with indices in the image
of~$\incO{O}$ is \m{\IncM{O}s_{w}^{\mathsf{T}}}; so these rows vanish in the columns outside the image, and the
$\incO{O}$-block of~$\conj{g}_{w}$ is~$s_{w}$. Hence \m{j=(j_{w})\in\LO{O}} with
\m{\etaO{O}(j)=\prod_{w\in\QO{O}}\psiO{O}{w}(\det s_{w})^{-1}}, so \m{g_{A}^{-1}g\in\KO{O}} and
\begin{align*}
\pi(g)\,\phi&=\pi(g_{A})\,\pi(g_{A}^{-1}g)\,\phi=\etaO{O}(j)\,\pi(g_{A})\,\phi ,\\
\pi(g_{A})\,\phi&=\prod_{w\in\MO{O}}\psiO{O}{w}(\det s_{w})^{-1}\,\chact{O}{\pi}{A}\,\phi ,
\end{align*}
the second equation by Setting~\ref{s:aux}(c).
\end{proof}
\leanline{\lean{smulchosenGLRepeq}}

\begin{lemma}[the conjugate scalar function of a form is automorphic]\label{l:automorphic}
Let $O$ be an auxiliary pair (Setting~\ref{s:aux}), $L$ a linear form on \m{\Vk{n}{\kap{O}}} and
\m{\Xi\in\FormO{O}}. Then the function \m{\conj{\coef{L}{\Xi}}} is automorphic on \m{\std{E}{n}}.
\end{lemma}
\begin{proof}
By (Q5), \m{\coef{L}{\Xi}} is automorphic, and the three defining conditions are stable under complex conjugation.
\end{proof}
\leanline{\lean{isAutomorphicstaralgFormCoeff}}

\sect{Proof of Theorem 1}{s:refl}
Throughout this section $O$, $\Xi$, $L$ are as in Theorem~1, \m{W=\WO{O}}, \m{V=\std{E}{n}}, \m{\chi=\chO{O}},
\m{\omega=\omega^{W,V}_{\chi,1}}, and \m{P=\plO{O}}. Functions at~$P$ are functions on the \m{n\times r} matrices
over~$E_{P}$, as in (Q3).

\subhead{Step 0: reductions.} Let $f$ be a datum with \m{\PS{S}f=f} and \m{u\in\Gad{\sgO{O}}}. By
Lemma~\ref{l:automorphic} and (Q1), the descent of \m{\conj{\coef{L}{\Xi}}} at~$u$ is a continuous linear functional of
the datum. By (Q2) it suffices to show that it vanishes on \m{\pvO{O}{p}\otimes\varrho(z)\vacf} for every
polynomial~$p$ and every finite vector~$z$ supported on~$P$ and its conjugates; by (Q3)(a), on
\m{\pvO{O}{p}\otimes\Psi\cdot\vacf} for every Schwartz--Bruhat function~$\Psi$ at~$P$; and by (Q3)(b), since the
descent is linear in~$\Psi$, for the functions \m{\Psi=\cosI{\mathfrak{t}}{z}}. So fix $p$ and an idele~$\mathfrak{t}$
with integral components, and put, for a matrix~$z$ over~$E_{P}$ and \m{\Xi'\in\FormO{O}},
\[
C_{z}(\Xi')=\Desc{W}{V}{\chi}{1}\bigl(\pvO{O}{p}\otimes\cosI{\mathfrak{t}}{z}\cdot\vacf,\ \conj{\coef{L}{\Xi'}}\bigr)(u).
\]
It is conjugate-linear in~$\Xi'$. We have to show \m{C_{z}(\Xi)=0} for every~$z$.

\subhead{Step 1: transport of the centre.}
\begin{lemma}\label{l:transport}
For \m{w\in P} let $A_{w}$,~$B_{w}$ be \m{n\times r} matrices over~$\oo_{w}$ and \m{s_{w}\in\mathrm{GL}_{r}(\FF_{w})}, and
let $z$,~$z'$ be matrices over~$E_{P}$. Suppose that for every \m{w\in P}: \textup{(i)} $s_{w}$ is the reduction
of \m{1+B_{w}^{\mathsf{T}}A_{w}}; \textup{(ii)} the reduction~$\sigma_{w}$ of the \m{n\times n} matrix
\m{1+A_{w}B_{w}^{\mathsf{T}}} satisfies \m{\sigma_{w}\IncM{O}=\IncM{O}s_{w}} if \m{w\in\MO{O}} and
\m{\sigma_{w}^{\mathsf{T}}\IncM{O}=\IncM{O}s_{w}^{\mathsf{T}}} if \m{w\in\QO{O}}; \textup{(iii)} the component at~$w$ of
\m{\mathfrak{t}^{-1}\bigl(z'-(1+A_{w}B_{w}^{\mathsf{T}})z\bigr)} has entries in~$\oo_{w}$. Then
\[
\prod_{w\in P}\psiO{O}{w}(\det s_{w})\cdot C_{z'}\bigl(\chact{O}{\mathrm{R}}{(s_{w})_{w\in\MO{O}}}\,\Xi\bigr)=C_{z}(\Xi),
\]
where \m{\chact{O}{\mathrm{R}}{\,\cdot\,}} is the action of Setting~\ref{s:aux}\textup{(c)} on \m{\FormO{O}}.
\end{lemma}
\begin{proof}
The determinant of \m{1+A_{w}B_{w}^{\mathsf{T}}} equals that of \m{1+B_{w}^{\mathsf{T}}A_{w}}, whose reduction is
\m{\det s_{w}\ne0}; so \m{1+A_{w}B_{w}^{\mathsf{T}}\in\mathrm{GL}_{n}(\oo_{w})} by (Q6). By Lemma~\ref{l:atplace} there is
\m{g\in\Gfin{n}}, supported on~$P$, with these components; it lies in~\m{\Kz{n}}, and \m{\red{w}g=\sigma_{w}}. By (Q1),
with \m{u'=1} and \m{h'=(1,g)},
\[
C_{z}(\Xi)=\Desc{W}{V}{\chi}{1}\Bigl(\omega\bigl(1,(1,g)\bigr)\bigl(\pvO{O}{p}\otimes\cosI{\mathfrak{t}}{z}\cdot\vacf\bigr),\ \conj{\coef{L}{\Xi}}(\,\cdot\,g)\Bigr)(u).
\]
By (Q4)(a) the datum is \m{\pvO{O}{p}\otimes\Psi'\cdot\vacf} with \m{\Psi'(X)=\cosI{\mathfrak{t}}{z}(g^{-1}X)}, the
characteristic function of \m{gz+\mathfrak{t}\operatorname{Box}}, since $g$ preserves \m{\operatorname{Box}}; by (iii)
this set is \m{z'+\mathfrak{t}\operatorname{Box}}. By (Q5), \m{\coef{L}{\Xi}(hg)=\coef{L}{(\mathrm{R}(g)\Xi)}(h)}, and by
Lemma~\ref{l:level} and (ii),
\[
\mathrm{R}(g)\,\Xi=\prod_{w\in P}\psiO{O}{w}(\det s_{w})^{-1}\cdot\chact{O}{\mathrm{R}}{(s_{w})_{w\in\MO{O}}}\,\Xi .
\]
As $C_{z'}$ is conjugate-linear and the values \m{\psiO{O}{w}(\det s_{w})} have absolute value~$1$, the scalar comes
out as \m{\prod_{w}\psiO{O}{w}(\det s_{w})}.
\end{proof}

\subhead{Step 2: annihilators.} Fix~$z$. For \m{w\in P} let $\delta_{w}$ be the diagonal \m{n\times n} matrix with
the entry~$\varpi_{w}$ at the indices outside the image of~$\incO{O}$ if \m{w\in\QO{O}}, and all other entries~$1$; so
\m{\delta_{w}=1} for \m{w\in\MO{O}}. Apply Lemma~\ref{l:annihilator} to the matrix with \m{n+r} rows
\[
X_{w}=\begin{pmatrix}\delta_{w}z_{w}\\ \mathfrak{t}_{w}1_{r}\end{pmatrix},
\]
with the injection of~$[r]$ into the first~$n$ rows by~$\incO{O}$; it has the left inverse
\m{(0,\ \mathfrak{t}_{w}^{-1}1_{r})}. We get a square matrix~$\ell_{w}$ of size \m{n+r} over~$\oo_{w}$ with
\m{\ell_{w}X_{w}=0}. For a row index~$i_{0}$ of~$\ell_{w}$ let \m{\conj{\ell}_{w}(i_{0})\in\FF_{w}^{r}} be the vector of the
residues of the entries \m{(\ell_{w})_{i_{0},\incO{O}(c)}}, \m{c\in[r]}. There are two cases.

\subhead{Step 3: all residues vanish.} Suppose \m{\conj{\ell}_{w}(i_{0})=0} for all~$w$ and~$i_{0}$. Lemma~\ref{l:annihilator}
gives \m{\alpha_{w}\in\mathrm{GL}_{r}(E_{w})} such that \m{\delta_{w}z_{w}\alpha_{w}-\IncM{O}} and \m{\mathfrak{t}_{w}\alpha_{w}} have
entries of absolute value \m{<1}. By Lemma~\ref{l:atplace} there is \m{\alpha\in\Gfin{r}}, supported on~$P$, with
the components~$\alpha_{w}$. Let~$\Psi$ be the function at~$P$ with
\m{\Psi(X)=\PsiO{O}{((X\alpha)_{w}\bmod w)_{w}}} if \m{X\alpha} has entries in \m{\prod_{w\in P}\oo_{w}}, and \m{\Psi(X)=0}
otherwise. We collect four facts.

(1) \m{\Psi(z)=1}. At \m{w\in\MO{O}}, \m{z_{w}\alpha_{w}} is congruent to~\m{\IncM{O}}. At \m{w\in\QO{O}}, the rows of
\m{z_{w}\alpha_{w}} in the image of~$\incO{O}$ are congruent to those of~\m{\IncM{O}}, and the other rows, multiplied
by~$\varpi_{w}$, have absolute value \m{<1}, so they are integral. Hence \m{z\alpha} is integral, reduces to~\m{\IncM{O}}
at~$\MO{O}$, and its $\incO{O}$-rows reduce to the identity at~$\QO{O}$.

(2) $\Psi$ is Schwartz--Bruhat and invariant under the translations by \m{\mathfrak{t}\operatorname{Box}}: for~$X'$
in this set, \m{X'\alpha} has entries of absolute value \m{<1}, and (Q3)(d) applies. By (Q3)(c),
\m{\Psi=\sum_{z'\in\mathcal{Z}}\Psi(z')\cosI{\mathfrak{t}}{z'}} with a finite set~$\mathcal{Z}$ and all \m{\Psi(z')\ne0}; by~(1),
$\Psi\ne0$, so $\mathcal{Z}$ is not empty.

(3) For every~$z'$ with \m{\Psi(z')\ne0}: \m{\Psi(z')\,C_{z'}(\Xi)=C_{\IncM{O}\alpha^{-1}}(\Xi)}. Indeed, \m{X=z'\alpha} is
integral, congruent to~\m{\IncM{O}} at~$\MO{O}$, and its $\incO{O}$-rows have invertible reductions~$s_{w}$ at
\m{w\in\QO{O}}; put \m{s_{w}=1} for \m{w\in\MO{O}}. Apply Lemma~\ref{l:transport} with \m{A_{w}=X_{w}-\IncM{O}},
\m{B_{w}=\IncM{O}}, the centres \m{\IncM{O}\alpha^{-1}} and~$z'$. Since \m{\IncM{O}^{\mathsf{T}}\IncM{O}=1}:
\m{1+B_{w}^{\mathsf{T}}A_{w}} is the matrix of the $\incO{O}$-rows of~$X_{w}$, which gives~(i);
\m{(1+A_{w}B_{w}^{\mathsf{T}})\IncM{O}=X_{w}}, whose reduction is \m{\IncM{O}} at~$\MO{O}$, and the rows of
\m{1+A_{w}B_{w}^{\mathsf{T}}} in the image of~$\incO{O}$ are those of \m{X_{w}\IncM{O}^{\mathsf{T}}}, which gives~(ii);
and \m{(1+A_{w}B_{w}^{\mathsf{T}})\IncM{O}\alpha_{w}^{-1}=X_{w}\alpha_{w}^{-1}=z'_{w}}, which gives~(iii). The lemma yields
\m{\prod_{w\in\QO{O}}\psiO{O}{w}(\det s_{w})\,C_{z'}(\Xi)=C_{\IncM{O}\alpha^{-1}}(\Xi)}, and the product is \m{\Psi(z')}.

(4) By hypothesis~(c) of the theorem at the point \m{u\alpha}, by (Q1) with \m{u'=(1,\alpha)}, \m{h'=1}, and by (Q4)(b)
applied to the component \m{X\mapsto\PsiO{O}{(X\bmod w)_{w}}} of~\m{\Ups{O}}:
\begin{align*}
0&=\Desc{W}{V}{\chi}{1}\bigl(\pvO{O}{p}\otimes\Ups{O},\,\conj{\coef{L}{\Xi}}\bigr)(u\alpha)\\
&=a\cdot\Desc{W}{V}{\chi}{1}\bigl(\pvO{O}{p}\otimes\Psi\cdot\vacf,\,\conj{\coef{L}{\Xi}}\bigr)(u)
=a\sum_{z'\in\mathcal{Z}}\Psi(z')\,C_{z'}(\Xi)
\end{align*}
with a complex number \m{a\ne0}. By~(3), applied to every \m{z'\in\mathcal{Z}} and to~$z$, each term is
\m{C_{\IncM{O}\alpha^{-1}}(\Xi)=\Psi(z)C_{z}(\Xi)=C_{z}(\Xi)}. So \m{\#\mathcal{Z}\cdot C_{z}(\Xi)=0}, and
\m{C_{z}(\Xi)=0}.

\subhead{Step 4: some residue does not vanish.} Suppose \m{\conj{\ell}_{w_{0}}(i_{0})\ne0} for some \m{w_{0}\in P} and
some row index~$i_{0}$. First, for every family \m{s=(s_{w})_{w\in P}} of invertible matrices of the form
\m{s_{w}=1+\xi_{w}\conj{\ell}_{w}(i_{0})^{\mathsf{T}}}, \m{\xi_{w}\in\FF_{w}^{r}}:
\begin{equation}\label{e:same}
\prod_{w\in P}\psiO{O}{w}(\det s_{w})\cdot C_{z}\bigl(\chact{O}{\mathrm{R}}{(s_{w})_{w\in\MO{O}}}\,\Xi\bigr)=C_{z}(\Xi).
\end{equation}
This is Lemma~\ref{l:transport} with \m{z'=z}, \m{A_{w}=\IncM{O}} and \m{B_{w}=\mu_{w}\tilde{\xi}_{w}^{\mathsf{T}}}, where
\m{\tilde{\xi}_{w}} is a lift of~$\xi_{w}$ to~$\oo_{w}^{r}$ and \m{\mu_{w}\in\oo_{w}^{n}} is the vector with the entries
\m{(\ell_{w})_{i_{0},i}(\delta_{w})_{ii}}, \m{i\in[n]}. Indeed the matrix
\[
1+B_{w}^{\mathsf{T}}A_{w}=1+\tilde{\xi}_{w}\,(\IncM{O}^{\mathsf{T}}\mu_{w})^{\mathsf{T}}
\]
reduces to~$s_{w}$, because $\delta_{w}$ is~$1$ at the indices in the image of~$\incO{O}$; this is~(i). The matrix
\m{1+A_{w}B_{w}^{\mathsf{T}}=1+\IncM{O}\tilde{\xi}_{w}\mu_{w}^{\mathsf{T}}} differs from~$1$ only in the rows of the
image of~$\incO{O}$; its $\incO{O}$-block reduces to~$s_{w}$, and its entries in these rows and in the columns
outside the image are multiples of~$\varpi_{w}$ when \m{w\in\QO{O}}; this gives~(ii) at both kinds of places.
Finally \m{\mu_{w}^{\mathsf{T}}z_{w}} is the part of the row~$i_{0}$ of~$\ell_{w}$ in the first $n$ columns times
\m{\delta_{w}z_{w}}, which by \m{\ell_{w}X_{w}=0} equals \m{-\mathfrak{t}_{w}} times the part of that row in the last
$r$ columns; so \m{\mathfrak{t}_{w}^{-1}(z_{w}-(1+A_{w}B_{w}^{\mathsf{T}})z_{w})=-\mathfrak{t}_{w}^{-1}\IncM{O}\tilde{\xi}_{w}\mu_{w}^{\mathsf{T}}z_{w}}
is integral; this is~(iii).

\emph{If \m{w_{0}\in\MO{O}}.} Sum \eqref{e:same} over the families~$s$ with \m{s_{w}=1} for \m{w\in\QO{O}} and $s_{w}$,
\m{w\in\MO{O}}, running over all invertible matrices of the form \m{1+\xi\,\conj{\ell}_{w}(i_{0})^{\mathsf{T}}}. By
Definition~\ref{d:refl} and conjugate-linearity, the left sides add up to the value of~$C_{z}$ at
\m{\bigotimes_{w\in\MO{O}}\Refl{\psiO{O}{w}}{\conj{\ell}_{w}(i_{0})}} applied to~$\Xi$, which is~$0$ by hypothesis~(b), the vector
at~$w_{0}$ being non-zero. The right sides add up to the number of families, which is at least~$1$,
times~$C_{z}(\Xi)$. So \m{C_{z}(\Xi)=0}.

\emph{If \m{w_{0}\in\QO{O}}.} By hypothesis~(a) there is \m{a\in\FF_{w_{0}}} with \m{a\ne0} and
\m{\psiO{O}{w_{0}}(a)\ne1}. Choose \m{c_{0}\in[r]} with the entry \m{\conj{\ell}_{w_{0}}(i_{0})_{c_{0}}\ne0}, and let
\m{\xi_{w_{0}}} be \m{(a-1)/\conj{\ell}_{w_{0}}(i_{0})_{c_{0}}} times the basis vector of index~$c_{0}$; then
\m{s_{w_{0}}=1+\xi_{w_{0}}\conj{\ell}_{w_{0}}(i_{0})^{\mathsf{T}}} has determinant~$a$. With \m{s_{w}=1} for \m{w\ne w_{0}},
\eqref{e:same} reads \m{\psiO{O}{w_{0}}(a)\,C_{z}(\Xi)=C_{z}(\Xi)}. So \m{C_{z}(\Xi)=0}.

This proves Theorem~1.
\leanline{\lean{descentsVanishofreflectionSum}}

\sect{A maximum principle}{s:max}
\begin{proposition}[bounded functions killed by a Hecke polynomial at a small argument]\label{p:max}
Let $\varepsilon'$ be a family of signs on~$[r']$, $F$ a complex function on \m{\Gad{\varepsilon'}}, $w$ a finite
place of~$E$ and $T$ a complex number. Suppose that $w$ is split, \m{w\ne\conj{w}}; that $F$ is spherical at~$w$
for \m{\Rt{\varepsilon'}}; that \m{\Hk{\Rt{\varepsilon'}}{w}{T}\,F=0}; that \m{q_{w}^{\,r'+1}\,|T|\le1}; and that $|F|$ is
bounded. Then \m{F=0}.
\end{proposition}
\begin{proof}
By (Q7), for every \m{u\in\Gad{\varepsilon'}},
\m{0=\sum_{U}(-T)^{\dim U}q_{w}^{\dim U(\dim U-1)/2}F(u\,g_{U})}. Replace $u$ by \m{u\,g_{0}^{-1}}, with $g_{0}$ the
element of the zero subspace: the term of \m{U=0} is~$F(u)$, so
\m{|F(u)|\le\sum_{U\ne0}|T|^{\dim U}q_{w}^{\dim U(\dim U-1)/2}\,\sup|F|}. By the second part of (Q7) at \m{-|T|}, the
sum of the coefficients over all~$U$ is \m{\prod_{0\le a<r'}(1+q_{w}^{a}|T|)\le\exp\bigl(|T|\sum_{a<r'}q_{w}^{a}\bigr)}. As
\m{q_{w}\ge2}, \m{\sum_{a<r'}q_{w}^{a}<q_{w}^{r'}}, so the exponent is less than \m{q_{w}^{r'}|T|\le q_{w}^{-1}\le\tfrac12},
which is less than \m{\log2}. Hence the sum of the coefficients over \m{U\ne0} is less than~$1$, and the
inequality forces \m{\sup|F|=0}.
\end{proof}
\leanline{\lean{eqzeroofheckePolyOpeqzero}\sep in the background library: \lnm{maximum\_principle}}

\sect{Proof of Theorem 2}{s:tower}
\begin{lemma}[the zero term has vanishing integral]\label{l:zero}
In Setting~\ref{s:weil}, let $W$, $W'$ and~$V$ be hermitian spaces with points, of signs $\varepsilon$ on~$[r]$,
$\varepsilon'$ on~$[r']$ and \m{\one{n}} on~$[n]$, $\chi$ an automorphic character, $S$ a set of finite places
and $Y$ a complex function on \m{\Gad{\one{n}}} with \m{\DV{W}{V}{\chi}{S}{Y}}. Let $\Pi$ be an adapted splitting
of~$W'$ over~$W$ with planes indexed by a finite set~$[\nu]$. Let \m{f\in\mathcal{S}(W',V)} with \m{\PS{S}f=f} and
\m{u\in\Gad{\varepsilon'}}. Then
\[
\int_{[V]}\theta_{W,V}\Bigl(\ZT{\Pi}{V}{0}\bigl(\omega^{W',V}_{\chi\upsilon_{1}^{\nu},1}(u,h)f\bigr)\Bigr)\,Y(h)\,\dd h=0 .
\]
\end{lemma}
\begin{proof}
With $g$, $u_{0}$, $a$ as in (Q10)(a) and \m{f_{1}=\ZT{\Pi}{V}{0}(\omega^{W',V}_{\chi\upsilon_{1}^{\nu},1}(g^{-1},1)f)}, the
integral is \m{a\int_{[V]}\theta_{W,V}(\omega^{W,V}_{\chi,1}(u_{0},h)f_{1})Y(h)\,\dd h=a\cdot\Desc{W}{V}{\chi}{1}(f_{1},Y)(u_{0})}. By
(Q10)(b), \m{\PS{S}f_{1}=f_{1}}, so this vanishes by \m{\DV{W}{V}{\chi}{S}{Y}}.
\end{proof}
\leanline{\lean{integralsourceTermeqzero}}

\begin{proof}[Proof of Theorem 2]
By induction on~$\nu$, for all $W'$ and~$\Pi$ at once. Write \m{\chi_{1}=\chi\upsilon_{1}^{\nu}} and
\m{\omega'=\omega^{W',V}_{\chi_{1},1}}; let $f$ be a datum with \m{\PS{S}f=f} and
\m{F=\Desc{W'}{V}{\chi_{1}}{1}(f,Y)}. We show \m{F=0}.

\emph{The case \m{\nu=0}.} By (Q10)(c), \m{F(u)} is the integral of Lemma~\ref{l:zero}, which vanishes.

\emph{The case \m{\nu\ge1}.} By (Q1) and (Q2) we may assume \m{f=v^{\varepsilon',\one{n}}_{p}\otimes\varrho(z)\vacf} with a
polynomial~$p$ and a finite vector~$z$ whose components at~$S$ and its conjugates vanish; such a datum is
unramified at~$S$, in particular at~\m{\{\conj{w}\}}.

\emph{Step 1: a Hecke relation for~$F$.} Let $\Lambda$ be the map that sends a datum~$f'$ to the linear map
\m{Y'\mapsto\Desc{W'}{V}{\chi_{1}}{1}(f',Y')} from the automorphic functions on~$V$ to the functions on
\m{\Gad{\varepsilon'}}. On such linear maps~$\Lambda'$ let \m{g\in\Gfin{r'}} act by right translation of the values,
and \m{h\in\Gfin{n}} by the inverse translation of the argument:
\[
(\pi(g)\Lambda')(Y')=\Lambda'(Y')(\,\cdot\,g),\qquad
(\rho(h)\Lambda')(Y')=\Lambda'\bigl(Y'(\,\cdot\,h^{-1})\bigr).
\]
These are commuting representations. By (Q1), applied with \m{u'=(1,g)}, \m{h'=(1,h)} and the function
\m{Y'(\,\cdot\,h^{-1})}, we have \m{\Lambda(\omega'((1,g),(1,h))f')=\pi(g)\rho(h)\Lambda f'} for every datum~$f'$, which is
the hypothesis of (Q9) on~$\Lambda$; and the pairing \m{(Y',\Lambda')\mapsto\Lambda'(Y')} is invariant under the pairs
\m{(\Rt{\one{n}}(h),\rho(h))}, as (Q8) requires. Moreover \m{r'=r+2\nu} and \m{2r+2\nu+3<n}, so \m{r'<n}: the
condition \m{r'\le n} of (Q9) holds. The base change of~$\chi_{1}$ is unramified at~\m{\conj{w}}, by hypothesis for~$\chi$ and by (Q9) for~$\upsilon_{1}$. So
(Q9), at the split place~\m{\conj{w}}, shows that \m{\Lambda f} is spherical for $\pi$ and~$\rho$ and gives, for
$\rho$ on the left and $\pi$ on the right,
\[
\bigl(\Lambda f\bigr)\bigl(\Hk{\Rt{\one{n}}}{w}{T}\,Y\bigr)=\prod_{0\le a<n-r'}\bigl(1-q_{w}^{a}A\,T\bigr)\cdot\Hk{\Rt{\varepsilon'}}{w}{T'}\,F,
\qquad T'=q_{w}^{\,n-r'}AB^{-1}T ;
\]
here the left side has been rewritten by the adjunction (Q8) for the invariant pairing above, which moves the
Hecke polynomial of~$\rho$ at~\m{\conj{w}} on linear maps to the Hecke polynomial of right translation at~$w$
on~$Y$; on the right, evaluation at~$Y$ commutes with~$\pi$, and \m{F=(\Lambda f)(Y)} is spherical at~$w$. The left side is~$0$ by hypothesis. On the right, \m{|q_{w}^{a}AT|\le q_{w}^{\,n-r'-1+r'/2}\,|T|=q_{w}^{(r-r')/2}=q_{w}^{-\nu}<1}
for \m{a<n-r'}, because \m{|T|=q_{w}^{-(n-1-r/2)}} and \m{r'=r+2\nu}; so no factor of the product vanishes, and
\m{\Hk{\Rt{\varepsilon'}}{w}{T'}F=0}.

\emph{Step 2: the size of~$T'$.} We have \m{|T'|=q_{w}^{(n-r')/2}|T|}, so
\[
q_{w}^{\,r'+1}\,|T'|=q_{w}^{\,r'+1+(n-r')/2-(n-1-r/2)}=q_{w}^{(2r+2\nu+4-n)/2}\le1,
\]
because \m{2r+2\nu+3<n}.

\emph{Step 3: $F$ is bounded.} By (Q10)(d) there are a space~$W_{1}$ with an adapted splitting over~$W$ with
\m{\nu-1} planes and an adapted splitting~$\Pi'$ of~$W'$ over~$W_{1}$ with one plane. The induction hypothesis,
applied to~$W_{1}$, gives \m{\DV{W_{1}}{V}{\chi\upsilon_{1}^{\nu-1}}{S}{Y}}: its hypotheses are those of the theorem,
with \m{2r+2(\nu-1)+3<n}. By Lemma~\ref{l:zero} for~$\Pi'$ and the character \m{\chi\upsilon_{1}^{\nu-1}},
\[
\int_{[V]}\theta_{W_{1},V}\Bigl(\ZT{\Pi'}{V}{0}\bigl(\omega'(u,h)f\bigr)\Bigr)\,Y(h)\,\dd h=0\qquad\text{for every }u.
\]
Hence, with
\m{\thd=\theta_{W',V}-\theta_{W_{1},V}\circ\ZT{\Pi'}{V}{0}} as in (Q11)(c),
\[
F(u)=\int_{[V]}\thd\bigl(\omega'(u,h)f\bigr)\,Y(h)\,\dd h\qquad\text{for every }u\in\Gad{\varepsilon'} .
\]
Fix~$u$. If \m{\operatorname{ht}(u)>2}, (Q11)(b) gives a non-zero isotropic vector of the lattice at~$u$ of length
at most \m{2/\operatorname{ht}(u)} and a rational point~$\gamma$ of~$W'$ that carries it to~\m{k_{\Pi'}}; then
\m{k_{\Pi'}} has this length in the lattice at~\m{\gamma u}. If \m{\operatorname{ht}(u)\le2}, put \m{\gamma=1}. In both
cases (Q11)(d) gives \m{|\thd(\omega'(\gamma u,h)f)|\le C_{0}} for all~$h$, with a constant~$C_{0}$ that does not depend
on~$u$. By (Q1), \m{F(u)=F(\gamma u)}, so \m{|F(u)|\le C_{0}\,\sup|Y|} times the measure of~$[V]$.

\emph{Step 4.} By Steps 1 to~3 and Proposition~\ref{p:max}, applied to $F$, $w$ and~$T'$, we get \m{F=0}.
\end{proof}
\leanline{\lean{descentsVanishaddplanes}}

\sect{Where the results are used}{s:use}
\emph{In the unit on the rows.} Let $O$ be an auxiliary pair with a parameter that fits it, as in the second
note, and let \m{\Xi\in\FormO{O}} be killed by the class-sum lift, by the reflection sums of hypothesis~(b), and by
the Hecke polynomial at one split place~$w$ outside the level at a point~$T$ of the absolute value of Theorem~2.
Let first $L$ be the hermitian product with the highest vector of the second note, so that \m{\coef{L}{\Xi}} is
the scalar function of its first theorem. By that theorem the descent of~$\conj{\coef{L}{\Xi}}$ with the test datum
at the highest vector vanishes; the edge property extends this to all polynomial vectors, and passing from the
characters of the pair to the characters \m{(\chO{O},1)} only multiplies descents by a character. This is
hypothesis~(c) of Theorem~1 for this one linear form. The unit on the rows needs (c) for the hermitian products
with all vectors of \m{\Vk{n}{\kap{O}}}, and obtains it in two steps. For the product with a translate of the highest
vector by an element of \m{\Ginf{\one{n}}}, the scalar function is a right translate of the former one, and (Q1)
moves the translation to the datum, which remains a polynomial vector tensor~\m{\Ups{O}}. And the translates of
the highest vector span \m{\Vk{n}{\kap{O}}}, while the descent of~$\conj{\coef{L}{\Xi}}$ is additive in~$L$. Then
Theorem~1 gives
\m{\DV{\WO{O}}{\std{E}{n}}{\chO{O}}{S}{\conj{\coef{L}{\Xi}}}} with $S$ the places outside the level, and Theorem~2, applied
to the sum of~\m{\WO{O}} and of $\nu$ hyperbolic planes with its tautological splitting, gives vanishing descents
to this sum for the character \m{\chO{O}\upsilon_{1}^{\nu}}, provided \m{2r+2\nu+3<n}. The unit on the rows deduces that
such forms have vanishing theta periods along a definite subgroup, which makes the test forms of the proof
orthogonal to them: forms that are not lifts but share one Hecke root with the lifts do not interfere.

\emph{In the unit on traces.} Lemma~\ref{l:level} is used on its own, to identify the action of elements of
\m{\Kz{n}} with block reductions on vectors of the level of~$O$.

\sect{Notes on the formal text}{s:notes}
\begin{enumerate}[label=\textup{\arabic*.},leftmargin=1.6em,itemsep=5pt]
\item \emph{The characters.} In this unit the Weil representation of the pair is taken with the second
character trivial: the statements use \m{\omega^{W,V}_{\chi,1}}, also in hypothesis~(c) of Theorem~1, whereas the
second note uses the pair of characters of~$O$. The definition of vanishing descents has only one character as
an argument.
\item \emph{Forms and their scalar functions.} In Theorem~1 the function whose descents vanish is the complex
conjugate of~\m{\coef{L}{\Xi}}, and the source of the descent is the definite group. The formal statement takes~$\Xi$
as an implicit argument and~$L$ as the first explicit one. Lemma~\ref{l:level} is stated in the formal text for an
arbitrary representation~$\pi$ of \m{\Gfin{n}} and a vector of the level, with the family of reductions of~$g$ as
a separate implicit argument tied to~$g$ by an equation; here the reductions are written \m{\conj{g}_{w}}.
\item \emph{Steps read at the level of their role.} In the proof of Theorem~2, Step~3 follows lines 80 to~107 of
the formal proof. The formal text applies the library's decay estimate with the exponent of the height equal
to~$0$; its hypothesis at points of height greater than~$2$ asks for a short isotropic vector in the lattice of
the indefinite space, carried to the frame of the one-plane splitting by a rational point (Witt's theorem), for
a basis of bounded length of the lattice of the definite space, for the invariance of the functional~$\thd$
along the frame, and for the inequality of (Q11)(c); the height of the definite space is~$1$, which makes the
bound uniform. The reader of the unit described these lines at the level of their role, and so does (Q11)(d),
which is a summary of the estimate in the case used and not a quotation of its statement. I read the statement
of the estimate and the lines that supply its hypotheses, and did not trace the constants. The second reader
pointed out that the argument which the first version of the dossier called a plane index in these lines is the
hypothesis that the height exceeds~$2$; the dossier now reads it so, and the text above follows this reading,
which agrees with the Lean text.
\item \emph{The reflection argument.} The first version of the dossier spoke, in Step~3 of the proof of
Theorem~1, of ``the number of coset terms''. In the Lean text the finite set~$\mathcal{Z}$ is a set of centres at
which~$\Psi$ does not vanish, the coefficient of the indicator at~$z'$ is~$\Psi(z')$, and the identity of fact~(3)
makes every term equal to \m{\Psi(z)C_{z}(\Xi)}; the set is not empty because \m{\Psi\ne0}. This is the second
reader's correction, which the dossier now follows, and it is what I find in lines 129 to~139. The matrix to which Lemma~\ref{l:annihilator} is applied in Step~2 has the lower
block \m{\mathfrak{t}_{w}1_{r}}; its left inverse does not depend on the upper block.
\item \emph{Automation.} Several steps of both long proofs are closed by the simplifier with long lists of
lemmas: the verification of the conditions (i) to~(iii) of Lemma~\ref{l:transport} in its three uses, fact~(1), and
the arithmetic of Steps 1 and~2 in the proof of Theorem~2. The computations given above are mine and are
complete; that the simplifier performs these computations was not replayed. In Proposition~\ref{p:max} the formal
proof applies a general maximum principle of the library with the constant~$2$ to the sum over all subspaces;
the proof above is the same argument with the term of the zero subspace taken out.
\item \emph{References for the quoted facts.} From the catalogue of the handout and the second reader's
comparison with the literature; where no place is named, the \emph{reference is to be confirmed}. (Q7), the
$q$-binomial theorem in the form of a sum over subspaces: R.~P.~Stanley, \emph{Enumerative Combinatorics},
vol.~1, second edition, Proposition~1.7.2 and equation~(1.87), with the proof over a finite field; relation ``the
theorem'' (the second reader's verdict; the catalogue attributes the identity to Gauss and Cauchy). (Q9): the catalogue names S.~Rallis, \emph{Langlands'
functoriality and the Weil representation}, Amer.~J.~Math.~104 (1982), R.~Howe, \emph{$\theta$-series and
invariant theory} (Corvallis 1979), and T.~Tamagawa, Ann.~of~Math.~77 (1963), places not confirmed, relation
\emph{special case} (split places, unramified characters, the unramified vacuum; an identity of Hecke
polynomials obtained by sums over subspaces of the residue space, without the Satake transform); the second
reader locates the correspondence of unramified parameters in A.~Mínguez, \emph{Correspondance de Howe
explicite : paires duales de type~II}, Theorem~1 and Corollary~6.3, and marks the relation as \emph{unresolved}
for the identity of operators and its normalisation. (Q10)(a), the Iwasawa decomposition: J.~Arthur, \emph{A trace
formula for reductive groups~I}, §1, p.~917; relation \emph{special case}. The zero term \m{\ZT{\Pi}{V}{0}} is an
operator that the library chooses from a proved existence statement; the second reader compares it with the
Ikeda map and the partial Fourier transform of W.~T.~Gan, Y.~Qiu and S.~Takeda, \emph{The regularized
Siegel--Weil formula}, §2.7, and marks the relation as \emph{unresolved}: the published evaluation in the mixed
model uses different spaces (the catalogue compares with Kudla 1994 and Harris--Kudla--Sweet 1996, §1, without a
confirmed place). (Q11)(b), Witt's theorem: U.~A.~First, \emph{Witt's extension theorem for quadratic spaces
over semiperfect rings}, J.~Pure Appl.~Algebra 219 (2015), Corollary~4.9; relation \emph{special case} (the
extension theorem over a field, evaluated on the hermitian matrix of the space). (Q11)(d): the catalogue compares with the decay of theta kernels on Siegel sets in S.~Kudla,
S.~Rallis, \emph{A regularized Siegel--Weil formula: the first term identity}, Ann.~of~Math.~140 (1994), place
not confirmed, and notes that the method of the formal text is different (small shifts along a short isotropic
vector). (Q1) to (Q6), (Q8), (Q11)(a) and~(c): no reference recorded.
\end{enumerate}

\subhead{How this note was written.} From the dossier of the unit: eight members and seven structured proofs.
The program check of the unit's dossier covered \UnitSteps{} leaf steps of \UnitProofs{} proofs, found the
statements at their lines and compared the quotations of the imported declarations with the source. I read both
Lean files of the unit in full and compared every statement with the statements above, argument by argument. In
the background library I read the definitions behind Definitions \ref{d:dv} to~\ref{d:split} and the statements of
the lemmas behind (Q1) to (Q11); I read none of their proofs. The second reader of the artifact read the unit
independently; her two corrections to the dossier's words are taken over (notes 3 and~4), and I checked both
against the Lean lines. The proofs of the three lemmas of §\ref{s:level}, of Proposition~\ref{p:max} and of
Lemma~\ref{l:zero} follow the structured proofs step by step. The proof of Theorem~1 follows the formal one in its
four steps; the verifications inside Steps 3 and~4 are written out by me. The proof of Theorem~2 follows the
formal one; its Step~3 is at the level described in note~3. The remarks after the statements and §\ref{s:use} are
mine; for §\ref{s:use} I read the statement and the proof of the consuming lemma in the unit on the rows.

\head{Appendix A. Concordance with the Lean artifact}
{\fontsize{8.6}{11.6}\selectfont\renewcommand{\leanfont}{\ttfamily\fontsize{7.7}{10}\selectfont\color{leancol}}
\setlength{\LTleft}{0pt}\setlength{\LTright}{\fill}\setlength{\tabcolsep}{0pt}
\begin{longtable}{@{}p{26mm}p{64mm}p{38mm}r@{}}
\textbf{In this note} & \textbf{Lean declaration} & \textbf{File} & \textbf{Lines}\\\midrule\endhead
\ConcordanceRows
\end{longtable}}

\subhead{The two theorems as they are written in the artifact.}
\quoted{descentsVanishofreflectionSum}
\quoted{descentsVanishaddplanes}
The second statement is to be read with the variable line at the head of its file, which introduces $W$, $W'$,
$V$, the hypothesis of vanishing descents for~$W$ and the splitting:
{\leanfont (h0 : DescentsVanish W V a S Y) (D : W'.Splitting W p)}.
\subhead{Legend.} In the formal text {\leanfont y} is the form~$\Xi$ and {\leanfont u} the linear form~$L$;
{\leanfont algFormCoeff 1 u y.1} is \m{\coef{L}{\Xi}} and {\leanfont star} is complex conjugation;
{\leanfont O.places} is \m{\plO{O}} and {\leanfont O.places.conjClosureᶜ} the set~$S$ of Theorem~1;
{\leanfont O.formAction} is the action of Setting~\ref{s:aux}(c) on forms and {\leanfont O.chosenGLRep π} the
action \m{\chact{O}{\pi}{\,\cdot\,}}; {\leanfont reflectionSum} and {\leanfont piTensorElem} are
\m{\Refl{\psi}{y}} and the tensor product; {\leanfont lift (pairThetaKer W V a 1).swap f Y} is the descent
\m{\Desc{W}{V}{\chi}{1}(f,Y)}, with {\leanfont a} the character~$\chi$; {\leanfont projOn \_ S} is \m{\PS{S}};
{\leanfont heckePolyOp (finRightRep ..) hw X} is \m{\Hk{\Rt{\varepsilon'}}{w}{T}}, with {\leanfont X} the
number~$T$; {\leanfont w.q} is~$q_{w}$; the index types {\leanfont m} and {\leanfont p} are $[r']$ and~$[\nu]$, the
signs {\leanfont s}, {\leanfont t} are $\varepsilon$,~$\varepsilon'$, the splitting {\leanfont D} is~$\Pi$, and
{\leanfont upsilonU1} is~$\upsilon_{1}$; {\leanfont D.sourceTerm V 0} is \m{\ZT{\Pi}{V}{0}}; the points
{\leanfont g} of the adelic group of the indefinite space are~$u$, and the elements {\leanfont k} of
\m{\Kz{n}} are~$g$.

\head{Notation}
\begingroup\fontsize{8.6}{10.9}\selectfont\raggedright
\setlength{\columnsep}{9mm}\setlength{\parskip}{0pt}
\begin{multicols}{2}
\nt{$E$, $\tau$, $w$, $\conj{w}$, $w_{0}$}{a CM field with a CM type; an infinite place; finite places and conjugates}
\nt{$E_{w}$, $\oo_{w}$, $\FF_{w}$, $q_{w}$, $\varpi_{w}$, $|\cdot|_{w}$}{completion, integers, residue field, its
number of elements, a uniformiser, the absolute value}
\nt{$\AAf$, $\ohat$, $E_{P}$, $P$}{finite adeles, their integers; the product of the~$E_{w}$ over a finite set~$P$
of places}
\nt{$[r]$, $[n]$, $[r']$, $[\nu]$; $r$, $n$, $r'$, $\nu$}{finite index sets; their numbers of elements (ranks of
$W$, $V$, $W'$; number of planes)}
\nt{$i$, $i_{0}$; $c$, $c'$, $c_{0}$; $a$}{row indices; indices in $[r]$; a natural number, and in the
proofs also a complex number or an element of a finite field}
\nt{$X^{*}$, $X^{\mathsf{T}}$, $\conj{z}$}{conjugate transpose, transpose, conjugation}
\nt{$\upsilon_{1}$, $\chi$, $\chi'$, $\chi_{1}$}{the splitting character on norm one; automorphic characters}
\nt{$\varepsilon$, $\varepsilon'$, $\one{n}$}{families of signs; the constant family~$1$}
\nt{$\Ginf{\varepsilon}$, $\Gfin{r}$, $\Gad{\varepsilon}$, $\Kz{r}$, $\red{w}$}{the adelic unitary groups, the integral
points, reduction modulo~$w$}
\nt{$u$, $u'$, $u_{0}$; $h$, $h'$; $\gamma$; $g$, $g'$, $g_{A}$, $g_{U}$, $g_{0}$}{points of the adelic groups of the
indefinite spaces; of the definite space; rational points; elements of the finite adelic groups, mostly
integral}
\nt{$W$, $W'$, $W_{1}$, $V$, $J$, $J'$, $\std{E}{n}$, $[V]$}{hermitian spaces with points, their matrices, the standard
definite space, the automorphic quotient}
\nt{$\mathcal{S}(W,V)$, $\varrho(z)$, $\omega^{W,V}_{\chi,\chi'}$, $\omega$, $\omega'$, $\theta_{W,V}$}{data, Heisenberg
operators, the Weil representation of the pair, the theta functional}
\nt{$f$, $f'$, $f_{1}$, $\varphi$, $\vacf$, $v$, $v^{\varepsilon,\varepsilon'}_{p}$, $p$}{data; the vacuum; a Fock vector;
the vector of a polynomial~$p$ in the variables \m{X_{\tau,i,c}}}
\nt{$\Desc{W}{V}{\chi}{\chi'}(f,Y)$, $Y$, $Y'$, $\dd h$}{the theta descent of a function on the definite group;
the invariant measure}
\nt{$S$, $\PS{S}$, $\DV{W}{V}{\chi}{S}{Y}$}{a set of finite places; the projector onto data unramified at~$S$;
vanishing of descents (Definition~\ref{d:dv})}
\nt{$O$, $\SigO{O}$, $\MO{O}$, $\QO{O}$, $\plO{O}$}{an auxiliary pair, its level system, chosen places, twist places,
all its places}
\nt{$\WO{O}$, $\sgO{O}$, $\incO{O}$, $\IncM{O}$, $\psiO{O}{w}$, $\chO{O}$}{its space, signs, embedding of index sets, the
matrix of the embedding, its residue characters, its character}
\nt{$\kap{O}$, $\eO{O}$, $\Vk{n}{\kap{O}}$, $\rhoW{\kap{O}}{\eO{O}}$}{the weight on the definite side, its space and
representation}
\nt{$\LO{O}$, $j$, $\etaO{O}$, $\KO{O}$, $\tetaO{O}$}{the level in the residue groups, an element, its character; the
level and its character}
\nt{$\FormO{O}$, $\Xi$, $\Xi'$, $L$, $\coef{L}{\Xi}$, $\mathrm{R}(g')$}{forms; a linear form; the scalar function; right
translation}
\nt{$\pi$, $\pi'$, $\rho$, $\mathcal{V}$, $\mathcal{V}'$, $\mathcal{V}^{O}$, $\phi$, $\phi'$}{representations of finite
adelic groups, their spaces, the vectors of the level of~$O$, vectors}
\nt{$\ChGL{\Sigma}{r}$, $A$, $\chact{O}{\pi}{A}$}{the product of the \m{\mathrm{GL}_{r}(\FF_{w})} over the chosen places,
an element, its action}
\nt{$\PsiO{O}{X}$, $\Ups{O}$, $\pvO{O}{p}$}{the function of residue matrices, the finite datum, the polynomial
vectors of~$O$}
\nt{$\FF$, $\psi$, $y$, $\xi$, $\Refl{\psi}{y}$, $G_{w}$, $a_{w}$}{a finite field, a multiplicative character,
vectors, the reflection sum; finite groups and elements of their group algebras}
\nt{$T$, $T'$, $T_{w,a}$, $\Hk{\pi}{w}{T}$, $\Rt{\varepsilon'}$}{complex numbers; Hecke operators; the Hecke
polynomial; right translation of functions}
\nt{$U$, $F$, $\alpha_{i}$}{a subspace of \m{\FF_{w}^{r'}}; a function on the adelic group of~$W'$; Satake parameters}
\nt{$\Pi$, $\Pi'$, $k_{\Pi}$, $l_{\Pi}$, $i_{\Pi}$, $\hyp{W}{\nu}$, $\ZT{\Pi}{V}{0}$}{splittings, their frames, the sum
with $\nu$ hyperbolic planes, the zero term}
\nt{$\Psi$, $\Psi'$, $\Psi''$, $\operatorname{Box}$, $\mathfrak{t}$, $\cosI{\mathfrak{t}}{z}$, $\mathcal{Z}$}{functions at
a finite set of places, the integral matrices, an idele, the indicator of a coset, a finite set of centres}
\nt{$z$, $z'$, $X$, $X'$, $X_{w}$, $Z$, $Z'$}{vectors and matrices over products of local fields ($z$ also an
idele in Setting~\ref{s:weil})}
\nt{$C_{z}(\Xi)$}{the descent with the indicator of the coset of~$z$ (Step~0 of §\ref{s:refl})}
\nt{$A_{w}$, $B_{w}$, $s_{w}$, $\sigma_{w}$, $\conj{g}_{w}$}{integral matrices; elements of \m{\mathrm{GL}_{r}(\FF_{w})};
reductions}
\nt{$\iota$, $\operatorname{Inc}_{\iota}$, $\ell$, $\ell_{w}$, $\conj{\ell}_{w}(i_{0})$}{an injection of index sets, its
matrix; annihilators; rows of residues}
\nt{$\alpha$, $\alpha_{w}$, $\beta$, $\beta'$, $\delta_{w}$, $\mu_{w}$, $\xi_{w}$, $\tilde{\xi}_{w}$}{matrices and
vectors in the proofs of Lemma~\ref{l:annihilator} and Theorem~1 ($\beta$ is a bilinear map in (Q8))}
\nt{$\Lambda$, $\Lambda'$, $\thd$, $A$, $B$}{linear maps on data; the difference of theta functionals; in (Q9),
$A$ and~$B$ are complex numbers}
\nt{$\operatorname{ht}(u)$, $C$, $C_{0}$}{the height; constants}
\end{multicols}
\endgroup

\end{document}
